% Options for packages loaded elsewhere
% Options for packages loaded elsewhere
\PassOptionsToPackage{unicode}{hyperref}
\PassOptionsToPackage{hyphens}{url}
\PassOptionsToPackage{dvipsnames,svgnames,x11names}{xcolor}
%
\documentclass[
  11pt,
  letterpaper,
  DIV=11,
  numbers=noendperiod]{scrartcl}
\usepackage{xcolor}
\usepackage[margin=1in]{geometry}
\usepackage{amsmath,amssymb}
\setcounter{secnumdepth}{5}
\usepackage{iftex}
\ifPDFTeX
  \usepackage[T1]{fontenc}
  \usepackage[utf8]{inputenc}
  \usepackage{textcomp} % provide euro and other symbols
\else % if luatex or xetex
  \usepackage{unicode-math} % this also loads fontspec
  \defaultfontfeatures{Scale=MatchLowercase}
  \defaultfontfeatures[\rmfamily]{Ligatures=TeX,Scale=1}
\fi
\usepackage{lmodern}
\ifPDFTeX\else
  % xetex/luatex font selection
\fi
% Use upquote if available, for straight quotes in verbatim environments
\IfFileExists{upquote.sty}{\usepackage{upquote}}{}
\IfFileExists{microtype.sty}{% use microtype if available
  \usepackage[]{microtype}
  \UseMicrotypeSet[protrusion]{basicmath} % disable protrusion for tt fonts
}{}
\makeatletter
\@ifundefined{KOMAClassName}{% if non-KOMA class
  \IfFileExists{parskip.sty}{%
    \usepackage{parskip}
  }{% else
    \setlength{\parindent}{0pt}
    \setlength{\parskip}{6pt plus 2pt minus 1pt}}
}{% if KOMA class
  \KOMAoptions{parskip=half}}
\makeatother
% Make \paragraph and \subparagraph free-standing
\makeatletter
\ifx\paragraph\undefined\else
  \let\oldparagraph\paragraph
  \renewcommand{\paragraph}{
    \@ifstar
      \xxxParagraphStar
      \xxxParagraphNoStar
  }
  \newcommand{\xxxParagraphStar}[1]{\oldparagraph*{#1}\mbox{}}
  \newcommand{\xxxParagraphNoStar}[1]{\oldparagraph{#1}\mbox{}}
\fi
\ifx\subparagraph\undefined\else
  \let\oldsubparagraph\subparagraph
  \renewcommand{\subparagraph}{
    \@ifstar
      \xxxSubParagraphStar
      \xxxSubParagraphNoStar
  }
  \newcommand{\xxxSubParagraphStar}[1]{\oldsubparagraph*{#1}\mbox{}}
  \newcommand{\xxxSubParagraphNoStar}[1]{\oldsubparagraph{#1}\mbox{}}
\fi
\makeatother

\usepackage{color}
\usepackage{fancyvrb}
\newcommand{\VerbBar}{|}
\newcommand{\VERB}{\Verb[commandchars=\\\{\}]}
\DefineVerbatimEnvironment{Highlighting}{Verbatim}{commandchars=\\\{\}}
% Add ',fontsize=\small' for more characters per line
\newenvironment{Shaded}{}{}
\newcommand{\AlertTok}[1]{\textcolor[rgb]{1.00,0.33,0.33}{\textbf{#1}}}
\newcommand{\AnnotationTok}[1]{\textcolor[rgb]{0.42,0.45,0.49}{#1}}
\newcommand{\AttributeTok}[1]{\textcolor[rgb]{0.84,0.23,0.29}{#1}}
\newcommand{\BaseNTok}[1]{\textcolor[rgb]{0.00,0.36,0.77}{#1}}
\newcommand{\BuiltInTok}[1]{\textcolor[rgb]{0.84,0.23,0.29}{#1}}
\newcommand{\CharTok}[1]{\textcolor[rgb]{0.01,0.18,0.38}{#1}}
\newcommand{\CommentTok}[1]{\textcolor[rgb]{0.42,0.45,0.49}{#1}}
\newcommand{\CommentVarTok}[1]{\textcolor[rgb]{0.42,0.45,0.49}{#1}}
\newcommand{\ConstantTok}[1]{\textcolor[rgb]{0.00,0.36,0.77}{#1}}
\newcommand{\ControlFlowTok}[1]{\textcolor[rgb]{0.84,0.23,0.29}{#1}}
\newcommand{\DataTypeTok}[1]{\textcolor[rgb]{0.84,0.23,0.29}{#1}}
\newcommand{\DecValTok}[1]{\textcolor[rgb]{0.00,0.36,0.77}{#1}}
\newcommand{\DocumentationTok}[1]{\textcolor[rgb]{0.42,0.45,0.49}{#1}}
\newcommand{\ErrorTok}[1]{\textcolor[rgb]{1.00,0.33,0.33}{\underline{#1}}}
\newcommand{\ExtensionTok}[1]{\textcolor[rgb]{0.84,0.23,0.29}{\textbf{#1}}}
\newcommand{\FloatTok}[1]{\textcolor[rgb]{0.00,0.36,0.77}{#1}}
\newcommand{\FunctionTok}[1]{\textcolor[rgb]{0.44,0.26,0.76}{#1}}
\newcommand{\ImportTok}[1]{\textcolor[rgb]{0.01,0.18,0.38}{#1}}
\newcommand{\InformationTok}[1]{\textcolor[rgb]{0.42,0.45,0.49}{#1}}
\newcommand{\KeywordTok}[1]{\textcolor[rgb]{0.84,0.23,0.29}{#1}}
\newcommand{\NormalTok}[1]{\textcolor[rgb]{0.14,0.16,0.18}{#1}}
\newcommand{\OperatorTok}[1]{\textcolor[rgb]{0.14,0.16,0.18}{#1}}
\newcommand{\OtherTok}[1]{\textcolor[rgb]{0.44,0.26,0.76}{#1}}
\newcommand{\PreprocessorTok}[1]{\textcolor[rgb]{0.84,0.23,0.29}{#1}}
\newcommand{\RegionMarkerTok}[1]{\textcolor[rgb]{0.42,0.45,0.49}{#1}}
\newcommand{\SpecialCharTok}[1]{\textcolor[rgb]{0.00,0.36,0.77}{#1}}
\newcommand{\SpecialStringTok}[1]{\textcolor[rgb]{0.01,0.18,0.38}{#1}}
\newcommand{\StringTok}[1]{\textcolor[rgb]{0.01,0.18,0.38}{#1}}
\newcommand{\VariableTok}[1]{\textcolor[rgb]{0.89,0.38,0.04}{#1}}
\newcommand{\VerbatimStringTok}[1]{\textcolor[rgb]{0.01,0.18,0.38}{#1}}
\newcommand{\WarningTok}[1]{\textcolor[rgb]{1.00,0.33,0.33}{#1}}

\usepackage{longtable,booktabs,array}
\usepackage{calc} % for calculating minipage widths
% Correct order of tables after \paragraph or \subparagraph
\usepackage{etoolbox}
\makeatletter
\patchcmd\longtable{\par}{\if@noskipsec\mbox{}\fi\par}{}{}
\makeatother
% Allow footnotes in longtable head/foot
\IfFileExists{footnotehyper.sty}{\usepackage{footnotehyper}}{\usepackage{footnote}}
\makesavenoteenv{longtable}
\usepackage{graphicx}
\makeatletter
\newsavebox\pandoc@box
\newcommand*\pandocbounded[1]{% scales image to fit in text height/width
  \sbox\pandoc@box{#1}%
  \Gscale@div\@tempa{\textheight}{\dimexpr\ht\pandoc@box+\dp\pandoc@box\relax}%
  \Gscale@div\@tempb{\linewidth}{\wd\pandoc@box}%
  \ifdim\@tempb\p@<\@tempa\p@\let\@tempa\@tempb\fi% select the smaller of both
  \ifdim\@tempa\p@<\p@\scalebox{\@tempa}{\usebox\pandoc@box}%
  \else\usebox{\pandoc@box}%
  \fi%
}
% Set default figure placement to htbp
\def\fps@figure{htbp}
\makeatother





\setlength{\emergencystretch}{3em} % prevent overfull lines

\providecommand{\tightlist}{%
  \setlength{\itemsep}{0pt}\setlength{\parskip}{0pt}}



 


\usepackage{booktabs}
\usepackage{float}
\usepackage{fancyhdr}
\usepackage{amsmath}
\pagestyle{fancy}
\fancyhead[L]{\textit{Migration Networks \& Shock Propagation}}
\fancyhead[R]{\thepage}
\fancyfoot[C]{}
\KOMAoption{captions}{tableheading}
\makeatletter
\@ifpackageloaded{caption}{}{\usepackage{caption}}
\AtBeginDocument{%
\ifdefined\contentsname
  \renewcommand*\contentsname{Table of contents}
\else
  \newcommand\contentsname{Table of contents}
\fi
\ifdefined\listfigurename
  \renewcommand*\listfigurename{List of Figures}
\else
  \newcommand\listfigurename{List of Figures}
\fi
\ifdefined\listtablename
  \renewcommand*\listtablename{List of Tables}
\else
  \newcommand\listtablename{List of Tables}
\fi
\ifdefined\figurename
  \renewcommand*\figurename{Figure}
\else
  \newcommand\figurename{Figure}
\fi
\ifdefined\tablename
  \renewcommand*\tablename{Table}
\else
  \newcommand\tablename{Table}
\fi
}
\@ifpackageloaded{float}{}{\usepackage{float}}
\floatstyle{ruled}
\@ifundefined{c@chapter}{\newfloat{codelisting}{h}{lop}}{\newfloat{codelisting}{h}{lop}[chapter]}
\floatname{codelisting}{Listing}
\newcommand*\listoflistings{\listof{codelisting}{List of Listings}}
\makeatother
\makeatletter
\makeatother
\makeatletter
\@ifpackageloaded{caption}{}{\usepackage{caption}}
\@ifpackageloaded{subcaption}{}{\usepackage{subcaption}}
\makeatother
\usepackage{bookmark}
\IfFileExists{xurl.sty}{\usepackage{xurl}}{} % add URL line breaks if available
\urlstyle{same}
\hypersetup{
  pdftitle={Migration Networks as Shock Transmission Infrastructure},
  pdfauthor={Michael Fehl \& Xenia Alemany},
  colorlinks=true,
  linkcolor={blue},
  filecolor={Maroon},
  citecolor={Blue},
  urlcolor={Blue},
  pdfcreator={LaTeX via pandoc}}


\title{Migration Networks as Shock Transmission Infrastructure}
\usepackage{etoolbox}
\makeatletter
\providecommand{\subtitle}[1]{% add subtitle to \maketitle
  \apptocmd{\@title}{\par {\large #1 \par}}{}{}
}
\makeatother
\subtitle{Community Structure, Centrality, and the Illinois Budget
Impasse}
\author{Michael Fehl \& Xenia Alemany}
\date{2026-03-30}
\begin{document}
\maketitle

\renewcommand*\contentsname{Table of contents}
{
\hypersetup{linkcolor=}
\setcounter{tocdepth}{3}
\tableofcontents
}

\subsection{Data Pipeline}\label{data-pipeline}

All derived objects are built once and saved to \texttt{R\_objects/}. On
subsequent runs, the code checks whether each file already exists and
skips the expensive rebuild if so. To force a full rebuild, delete the
relevant file from \texttt{R\_objects/}.

\subsubsection{Cleaning helper}\label{cleaning-helper}

\begin{Shaded}
\begin{Highlighting}[]
\CommentTok{\# Standardize a raw MCAS Excel file: drop "Total" rows, rename columns,}
\CommentTok{\# forward{-}fill the Mexican state column, and add a US destination label.}
\NormalTok{clean\_dfs }\OtherTok{\textless{}{-}} \ControlFlowTok{function}\NormalTok{(df, state, }\AttributeTok{drop =} \ConstantTok{FALSE}\NormalTok{) \{}
\NormalTok{  dirty }\OtherTok{\textless{}{-}}\NormalTok{ df}
  \CommentTok{\# Drop administrative "Total" subtotal rows}
\NormalTok{  dirty }\OtherTok{\textless{}{-}}\NormalTok{ dirty }\SpecialCharTok{|\textgreater{}}
    \FunctionTok{filter}\NormalTok{(}\SpecialCharTok{!}\FunctionTok{str\_detect}\NormalTok{(}\FunctionTok{as.character}\NormalTok{(mx\_state),}
                       \FunctionTok{regex}\NormalTok{(}\StringTok{"Total"}\NormalTok{, }\AttributeTok{ignore\_case =} \ConstantTok{TRUE}\NormalTok{))) }\SpecialCharTok{|\textgreater{}}
    \FunctionTok{filter}\NormalTok{(}\SpecialCharTok{!}\FunctionTok{str\_detect}\NormalTok{(}\FunctionTok{as.character}\NormalTok{(mx\_municipality),}
                       \FunctionTok{regex}\NormalTok{(}\StringTok{"Total"}\NormalTok{, }\AttributeTok{ignore\_case =} \ConstantTok{TRUE}\NormalTok{)))}
  \CommentTok{\# Rename to standard column names}
\NormalTok{  dirty }\OtherTok{\textless{}{-}}\NormalTok{ dirty }\SpecialCharTok{|\textgreater{}}
    \FunctionTok{rename}\NormalTok{(}\StringTok{\textasciigrave{}}\AttributeTok{origin state}\StringTok{\textasciigrave{}} \OtherTok{=}\NormalTok{ mx\_state,}
           \StringTok{\textasciigrave{}}\AttributeTok{origin muni}\StringTok{\textasciigrave{}}  \OtherTok{=}\NormalTok{ mx\_municipality,}
           \AttributeTok{weight         =}\NormalTok{ pct\_matriculas)}
  \CommentTok{\# Forward{-}fill the state header rows that the raw Excel leaves blank}
\NormalTok{  dirty }\OtherTok{\textless{}{-}}\NormalTok{ dirty }\SpecialCharTok{|\textgreater{}} \FunctionTok{fill}\NormalTok{(}\StringTok{\textasciigrave{}}\AttributeTok{origin state}\StringTok{\textasciigrave{}}\NormalTok{, }\AttributeTok{.direction =} \StringTok{"down"}\NormalTok{)}
  \CommentTok{\# Add US destination label ("New\_York" {-}\textgreater{} "New York")}
\NormalTok{  dirty }\OtherTok{\textless{}{-}}\NormalTok{ dirty }\SpecialCharTok{|\textgreater{}} \FunctionTok{mutate}\NormalTok{(}\AttributeTok{destination =} \FunctionTok{str\_replace\_all}\NormalTok{(state, }\StringTok{"\_"}\NormalTok{, }\StringTok{" "}\NormalTok{))}
  \ControlFlowTok{if}\NormalTok{ (drop) dirty }\OtherTok{\textless{}{-}}\NormalTok{ dirty }\SpecialCharTok{|\textgreater{}}
    \FunctionTok{select}\NormalTok{(}\StringTok{\textasciigrave{}}\AttributeTok{origin state}\StringTok{\textasciigrave{}}\NormalTok{, }\StringTok{\textasciigrave{}}\AttributeTok{origin muni}\StringTok{\textasciigrave{}}\NormalTok{, destination, weight)}
\NormalTok{  dirty}
\NormalTok{\}}
\end{Highlighting}
\end{Shaded}

\subsubsection{\texorpdfstring{\texttt{long\_panel}: long-format
migration panel (CA, NY, IL ×
2010--2024)}{long\_panel: long-format migration panel (CA, NY, IL × 2010--2024)}}\label{long_panel-long-format-migration-panel-ca-ny-il-20102024}

\begin{Shaded}
\begin{Highlighting}[]
\ControlFlowTok{if}\NormalTok{ (}\FunctionTok{file.exists}\NormalTok{(}\FunctionTok{paste0}\NormalTok{(out\_dir, }\StringTok{"long\_panel.csv"}\NormalTok{))) \{}

  \CommentTok{\# {-}{-}{-} Fast path: load from disk {-}{-}{-}{-}{-}{-}{-}{-}{-}{-}{-}{-}{-}{-}{-}{-}{-}{-}{-}{-}{-}{-}{-}{-}{-}{-}{-}{-}{-}{-}{-}{-}{-}{-}{-}{-}{-}{-}{-}{-}{-}{-}{-}{-}}
\NormalTok{  long\_panel }\OtherTok{\textless{}{-}} \FunctionTok{read.csv}\NormalTok{(}\FunctionTok{paste0}\NormalTok{(out\_dir, }\StringTok{"long\_panel.csv"}\NormalTok{),}
                         \AttributeTok{stringsAsFactors =} \ConstantTok{FALSE}\NormalTok{) }\SpecialCharTok{|\textgreater{}}
    \FunctionTok{mutate}\NormalTok{(}\AttributeTok{year =} \FunctionTok{as.integer}\NormalTok{(year))}
  \FunctionTok{cat}\NormalTok{(}\StringTok{"Loaded long\_panel from disk:"}\NormalTok{, }\FunctionTok{nrow}\NormalTok{(long\_panel), }\StringTok{"rows}\SpecialCharTok{\textbackslash{}n}\StringTok{"}\NormalTok{)}

\NormalTok{\} }\ControlFlowTok{else}\NormalTok{ \{}

  \CommentTok{\# {-}{-}{-} Build path: load, clean, aggregate, merge {-}{-}{-}{-}{-}{-}{-}{-}{-}{-}{-}{-}{-}{-}{-}{-}{-}{-}{-}{-}{-}{-}{-}{-}{-}{-}{-}{-}}
\NormalTok{  years  }\OtherTok{\textless{}{-}} \DecValTok{2010}\SpecialCharTok{:}\DecValTok{2024}
\NormalTok{  states }\OtherTok{\textless{}{-}} \FunctionTok{c}\NormalTok{(}\StringTok{"California"}\NormalTok{, }\StringTok{"New\_York"}\NormalTok{, }\StringTok{"Illinois"}\NormalTok{)}

\NormalTok{  df\_list }\OtherTok{\textless{}{-}} \FunctionTok{list}\NormalTok{()}
  \ControlFlowTok{for}\NormalTok{ (state\_i }\ControlFlowTok{in}\NormalTok{ states) \{}
    \ControlFlowTok{for}\NormalTok{ (yr }\ControlFlowTok{in}\NormalTok{ years) \{}
\NormalTok{      path  }\OtherTok{\textless{}{-}} \FunctionTok{file.path}\NormalTok{(us\_states\_path,}
                         \FunctionTok{paste0}\NormalTok{(}\StringTok{"Edos\_USA\_"}\NormalTok{, yr),}
                         \FunctionTok{paste0}\NormalTok{(}\FunctionTok{toupper}\NormalTok{(state\_i), }\StringTok{"\_"}\NormalTok{, yr, }\StringTok{".xlsx"}\NormalTok{))}
\NormalTok{      dirty }\OtherTok{\textless{}{-}} \FunctionTok{read\_excel}\NormalTok{(path)}
\NormalTok{      clean }\OtherTok{\textless{}{-}} \FunctionTok{clean\_dfs}\NormalTok{(}\AttributeTok{df =}\NormalTok{ dirty, }\AttributeTok{state =}\NormalTok{ state\_i, }\AttributeTok{drop =} \ConstantTok{TRUE}\NormalTok{)}
\NormalTok{      df\_list[[}\FunctionTok{paste0}\NormalTok{(}\FunctionTok{tolower}\NormalTok{(state\_i), }\StringTok{"\_"}\NormalTok{, yr)]] }\OtherTok{\textless{}{-}}\NormalTok{ clean}
\NormalTok{    \}}
\NormalTok{  \}}

  \CommentTok{\# Rename weight {-}\textgreater{} key, THEN aggregate to one row per municipality before joining.}
  \CommentTok{\# Without the group\_by + summarise step, duplicate sub{-}rows in the raw Excel files}
  \CommentTok{\# cause a Cartesian product explosion when full\_join is called 45 times.}
\NormalTok{  wide\_df }\OtherTok{\textless{}{-}} \FunctionTok{imap}\NormalTok{(df\_list, }\ControlFlowTok{function}\NormalTok{(df, key) \{}
\NormalTok{    df }\SpecialCharTok{|\textgreater{}}
      \FunctionTok{select}\NormalTok{(}\StringTok{\textasciigrave{}}\AttributeTok{origin state}\StringTok{\textasciigrave{}}\NormalTok{, }\StringTok{\textasciigrave{}}\AttributeTok{origin muni}\StringTok{\textasciigrave{}}\NormalTok{, weight) }\SpecialCharTok{|\textgreater{}}
      \FunctionTok{group\_by}\NormalTok{(}\StringTok{\textasciigrave{}}\AttributeTok{origin state}\StringTok{\textasciigrave{}}\NormalTok{, }\StringTok{\textasciigrave{}}\AttributeTok{origin muni}\StringTok{\textasciigrave{}}\NormalTok{) }\SpecialCharTok{|\textgreater{}}
      \FunctionTok{summarise}\NormalTok{(}\SpecialCharTok{!!}\AttributeTok{key :=} \FunctionTok{sum}\NormalTok{(weight, }\AttributeTok{na.rm =} \ConstantTok{TRUE}\NormalTok{), }\AttributeTok{.groups =} \StringTok{"drop"}\NormalTok{)}
\NormalTok{  \}) }\SpecialCharTok{|\textgreater{}}
    \FunctionTok{reduce}\NormalTok{(}\ControlFlowTok{function}\NormalTok{(a, b) \{}
      \FunctionTok{full\_join}\NormalTok{(a, b, }\AttributeTok{by =} \FunctionTok{c}\NormalTok{(}\StringTok{"origin state"}\NormalTok{, }\StringTok{"origin muni"}\NormalTok{))}
\NormalTok{    \}) }\SpecialCharTok{|\textgreater{}}
    \FunctionTok{mutate}\NormalTok{(}\FunctionTok{across}\NormalTok{(}\FunctionTok{where}\NormalTok{(is.numeric), }\SpecialCharTok{\textasciitilde{}}\FunctionTok{replace\_na}\NormalTok{(., }\DecValTok{0}\NormalTok{)))}

  \CommentTok{\# Melt to long format; split "california\_2016" {-}\textgreater{} state + year}
\NormalTok{  long\_panel }\OtherTok{\textless{}{-}}\NormalTok{ wide\_df }\SpecialCharTok{|\textgreater{}}
    \FunctionTok{pivot\_longer}\NormalTok{(}\AttributeTok{cols      =} \SpecialCharTok{{-}}\FunctionTok{c}\NormalTok{(}\StringTok{\textasciigrave{}}\AttributeTok{origin state}\StringTok{\textasciigrave{}}\NormalTok{, }\StringTok{\textasciigrave{}}\AttributeTok{origin muni}\StringTok{\textasciigrave{}}\NormalTok{),}
                 \AttributeTok{names\_to  =} \StringTok{"state\_year"}\NormalTok{,}
                 \AttributeTok{values\_to =} \StringTok{"weight"}\NormalTok{) }\SpecialCharTok{|\textgreater{}}
    \FunctionTok{separate}\NormalTok{(state\_year, }\AttributeTok{into =} \FunctionTok{c}\NormalTok{(}\StringTok{"state"}\NormalTok{, }\StringTok{"year"}\NormalTok{),}
             \AttributeTok{sep =} \StringTok{"\_(?=}\SpecialCharTok{\textbackslash{}\textbackslash{}}\StringTok{d\{4\}$)"}\NormalTok{) }\SpecialCharTok{|\textgreater{}}
    \FunctionTok{mutate}\NormalTok{(}\AttributeTok{year =} \FunctionTok{as.integer}\NormalTok{(year))}

  \FunctionTok{write.csv}\NormalTok{(long\_panel, }\FunctionTok{paste0}\NormalTok{(out\_dir, }\StringTok{"long\_panel.csv"}\NormalTok{),}
            \AttributeTok{row.names =} \ConstantTok{FALSE}\NormalTok{)}
  \FunctionTok{cat}\NormalTok{(}\StringTok{"Built and saved long\_panel:"}\NormalTok{, }\FunctionTok{nrow}\NormalTok{(long\_panel), }\StringTok{"rows}\SpecialCharTok{\textbackslash{}n}\StringTok{"}\NormalTok{)}
\NormalTok{\}}
\end{Highlighting}
\end{Shaded}

\subsubsection{\texorpdfstring{\texttt{r\_matrix\_idx}: average
migration weighting matrix
(wide)}{r\_matrix\_idx: average migration weighting matrix (wide)}}\label{r_matrix_idx-average-migration-weighting-matrix-wide}

\begin{Shaded}
\begin{Highlighting}[]
\ControlFlowTok{if}\NormalTok{ (}\FunctionTok{file.exists}\NormalTok{(}\FunctionTok{paste0}\NormalTok{(out\_dir, }\StringTok{"r\_matrix\_idx.csv"}\NormalTok{))) \{}

\NormalTok{  r\_matrix\_idx }\OtherTok{\textless{}{-}} \FunctionTok{read.csv}\NormalTok{(}\FunctionTok{paste0}\NormalTok{(out\_dir, }\StringTok{"r\_matrix\_idx.csv"}\NormalTok{),}
                            \AttributeTok{check.names =} \ConstantTok{FALSE}\NormalTok{,}
                            \AttributeTok{stringsAsFactors =} \ConstantTok{FALSE}\NormalTok{)}
  \FunctionTok{cat}\NormalTok{(}\StringTok{"Loaded r\_matrix\_idx from disk:"}\NormalTok{, }\FunctionTok{nrow}\NormalTok{(r\_matrix\_idx), }\StringTok{"rows x"}\NormalTok{,}
      \FunctionTok{ncol}\NormalTok{(r\_matrix\_idx), }\StringTok{"cols}\SpecialCharTok{\textbackslash{}n}\StringTok{"}\NormalTok{)}

\NormalTok{\} }\ControlFlowTok{else}\NormalTok{ \{}

\NormalTok{  r\_matrix\_idx }\OtherTok{\textless{}{-}} \FunctionTok{read\_excel}\NormalTok{(}
    \StringTok{"1\_network\_estimation/2\_migration\_matrix\_estimation/migration\_weighting\_matrices\_2/AVG\_WEIGHTING\_MATRIX.xlsx"}
\NormalTok{  ) }\SpecialCharTok{|\textgreater{}}
    \FunctionTok{select}\NormalTok{(mx\_state, mx\_municipality, }\FunctionTok{everything}\NormalTok{())}

  \FunctionTok{write.csv}\NormalTok{(r\_matrix\_idx, }\FunctionTok{paste0}\NormalTok{(out\_dir, }\StringTok{"r\_matrix\_idx.csv"}\NormalTok{),}
            \AttributeTok{row.names =} \ConstantTok{FALSE}\NormalTok{)}
  \FunctionTok{cat}\NormalTok{(}\StringTok{"Loaded and saved r\_matrix\_idx:"}\NormalTok{, }\FunctionTok{nrow}\NormalTok{(r\_matrix\_idx), }\StringTok{"rows}\SpecialCharTok{\textbackslash{}n}\StringTok{"}\NormalTok{)}
\NormalTok{\}}

\CommentTok{\# Extract the numeric matrix once for use throughout the paper}
\NormalTok{numeric\_cols }\OtherTok{\textless{}{-}} \FunctionTok{setdiff}\NormalTok{(}\FunctionTok{names}\NormalTok{(r\_matrix\_idx), }\FunctionTok{c}\NormalTok{(}\StringTok{"mx\_state"}\NormalTok{, }\StringTok{"mx\_municipality"}\NormalTok{))}
\NormalTok{W\_raw        }\OtherTok{\textless{}{-}} \FunctionTok{as.matrix}\NormalTok{(r\_matrix\_idx[, numeric\_cols])}
\FunctionTok{rownames}\NormalTok{(W\_raw) }\OtherTok{\textless{}{-}}\NormalTok{ r\_matrix\_idx}\SpecialCharTok{$}\NormalTok{mx\_municipality}
\end{Highlighting}
\end{Shaded}

\subsubsection{\texorpdfstring{\texttt{ds2\_pivot}: quarterly US-state
remittance outflows
(wide)}{ds2\_pivot: quarterly US-state remittance outflows (wide)}}\label{ds2_pivot-quarterly-us-state-remittance-outflows-wide}

\begin{Shaded}
\begin{Highlighting}[]
\ControlFlowTok{if}\NormalTok{ (}\FunctionTok{file.exists}\NormalTok{(}\FunctionTok{paste0}\NormalTok{(out\_dir, }\StringTok{"ds2\_pivot.csv"}\NormalTok{))) \{}

\NormalTok{  ds2\_pivot }\OtherTok{\textless{}{-}} \FunctionTok{read.csv}\NormalTok{(}\FunctionTok{paste0}\NormalTok{(out\_dir, }\StringTok{"ds2\_pivot.csv"}\NormalTok{),}
                         \AttributeTok{check.names =} \ConstantTok{FALSE}\NormalTok{,}
                         \AttributeTok{stringsAsFactors =} \ConstantTok{FALSE}\NormalTok{) }\SpecialCharTok{|\textgreater{}}
    \FunctionTok{mutate}\NormalTok{(}\AttributeTok{date =} \FunctionTok{as.Date}\NormalTok{(date))}
  \FunctionTok{cat}\NormalTok{(}\StringTok{"Loaded ds2\_pivot from disk:"}\NormalTok{, }\FunctionTok{nrow}\NormalTok{(ds2\_pivot), }\StringTok{"rows}\SpecialCharTok{\textbackslash{}n}\StringTok{"}\NormalTok{)}

\NormalTok{\} }\ControlFlowTok{else}\NormalTok{ \{}

  \CommentTok{\# Load pre{-}cleaned CSV (exported from the Banxico Excel source)}
\NormalTok{  ds2 }\OtherTok{\textless{}{-}} \FunctionTok{read.csv}\NormalTok{(}\StringTok{"R\_objects/remit\_us\_states.csv"}\NormalTok{,}
                  \AttributeTok{stringsAsFactors =} \ConstantTok{FALSE}\NormalTok{)}

\NormalTok{  ds2\_pivot }\OtherTok{\textless{}{-}}\NormalTok{ ds2 }\SpecialCharTok{|\textgreater{}}
    \FunctionTok{mutate}\NormalTok{(}
      \AttributeTok{month =}\NormalTok{ (quarter }\SpecialCharTok{{-}} \DecValTok{1}\NormalTok{L) }\SpecialCharTok{*} \DecValTok{3}\NormalTok{L }\SpecialCharTok{+} \DecValTok{1}\NormalTok{L,}
      \AttributeTok{date  =} \FunctionTok{as.Date}\NormalTok{(}\FunctionTok{sprintf}\NormalTok{(}\StringTok{"\%d{-}\%02d{-}01"}\NormalTok{, year, month))}
\NormalTok{    ) }\SpecialCharTok{|\textgreater{}}
    \FunctionTok{pivot\_wider}\NormalTok{(}\AttributeTok{id\_cols =}\NormalTok{ date, }\AttributeTok{names\_from =}\NormalTok{ us\_state,}
                \AttributeTok{values\_from =}\NormalTok{ remit\_musd, }\AttributeTok{values\_fill =} \DecValTok{0}\NormalTok{) }\SpecialCharTok{|\textgreater{}}
    \FunctionTok{arrange}\NormalTok{(date)}

  \FunctionTok{write.csv}\NormalTok{(ds2\_pivot, }\FunctionTok{paste0}\NormalTok{(out\_dir, }\StringTok{"ds2\_pivot.csv"}\NormalTok{),}
            \AttributeTok{row.names =} \ConstantTok{FALSE}\NormalTok{)}
  \FunctionTok{cat}\NormalTok{(}\StringTok{"Built and saved ds2\_pivot:"}\NormalTok{, }\FunctionTok{nrow}\NormalTok{(ds2\_pivot), }\StringTok{"rows}\SpecialCharTok{\textbackslash{}n}\StringTok{"}\NormalTok{)}
\NormalTok{\}}
\end{Highlighting}
\end{Shaded}

\subsubsection{\texorpdfstring{\texttt{R\_hat\_df}: imputed
municipality-level remittances (wide,
quarterly)}{R\_hat\_df: imputed municipality-level remittances (wide, quarterly)}}\label{r_hat_df-imputed-municipality-level-remittances-wide-quarterly}

\begin{Shaded}
\begin{Highlighting}[]
\ControlFlowTok{if}\NormalTok{ (}\FunctionTok{file.exists}\NormalTok{(}\FunctionTok{paste0}\NormalTok{(out\_dir, }\StringTok{"R\_hat\_df.csv"}\NormalTok{))) \{}

\NormalTok{  R\_hat\_df }\OtherTok{\textless{}{-}} \FunctionTok{read.csv}\NormalTok{(}\FunctionTok{paste0}\NormalTok{(out\_dir, }\StringTok{"R\_hat\_df.csv"}\NormalTok{),}
                        \AttributeTok{check.names =} \ConstantTok{FALSE}\NormalTok{,}
                        \AttributeTok{stringsAsFactors =} \ConstantTok{FALSE}\NormalTok{)}
  \FunctionTok{cat}\NormalTok{(}\StringTok{"Loaded R\_hat\_df from disk:"}\NormalTok{, }\FunctionTok{nrow}\NormalTok{(R\_hat\_df), }\StringTok{"rows x"}\NormalTok{,}
      \FunctionTok{ncol}\NormalTok{(R\_hat\_df), }\StringTok{"cols}\SpecialCharTok{\textbackslash{}n}\StringTok{"}\NormalTok{)}

\NormalTok{\} }\ControlFlowTok{else}\NormalTok{ \{}

  \CommentTok{\# Column{-}normalise W (each state column sums to 1)}
\NormalTok{  W\_col\_raw }\OtherTok{\textless{}{-}} \FunctionTok{sweep}\NormalTok{(W\_raw, }\DecValTok{2}\NormalTok{, }\FunctionTok{colSums}\NormalTok{(W\_raw), }\StringTok{"/"}\NormalTok{)}

  \CommentTok{\# Align states present in both W and DS2}
\NormalTok{  state\_cols\_ds2 }\OtherTok{\textless{}{-}} \FunctionTok{setdiff}\NormalTok{(}\FunctionTok{names}\NormalTok{(ds2\_pivot), }\StringTok{"date"}\NormalTok{)}
\NormalTok{  states\_common  }\OtherTok{\textless{}{-}} \FunctionTok{intersect}\NormalTok{(}\FunctionTok{colnames}\NormalTok{(W\_col\_raw), state\_cols\_ds2)}

\NormalTok{  W\_mat   }\OtherTok{\textless{}{-}}\NormalTok{ W\_col\_raw[, states\_common]}
\NormalTok{  DS2\_mat }\OtherTok{\textless{}{-}} \FunctionTok{as.matrix}\NormalTok{(ds2\_pivot[, states\_common])}

  \CommentTok{\# R\_hat = W \%*\% DS2\textquotesingle{}: (n\_munis x n\_states) \%*\% (n\_states x n\_quarters)}
\NormalTok{  R\_hat }\OtherTok{\textless{}{-}}\NormalTok{ W\_mat }\SpecialCharTok{\%*\%} \FunctionTok{t}\NormalTok{(DS2\_mat)}

\NormalTok{  R\_hat\_df }\OtherTok{\textless{}{-}}\NormalTok{ r\_matrix\_idx }\SpecialCharTok{|\textgreater{}}
    \FunctionTok{select}\NormalTok{(mx\_state, mx\_municipality) }\SpecialCharTok{|\textgreater{}}
    \FunctionTok{bind\_cols}\NormalTok{(}\FunctionTok{as.data.frame}\NormalTok{(R\_hat) }\SpecialCharTok{|\textgreater{}}
                \FunctionTok{setNames}\NormalTok{(}\FunctionTok{as.character}\NormalTok{(ds2\_pivot}\SpecialCharTok{$}\NormalTok{date)))}

  \FunctionTok{write.csv}\NormalTok{(R\_hat\_df, }\FunctionTok{paste0}\NormalTok{(out\_dir, }\StringTok{"R\_hat\_df.csv"}\NormalTok{),}
            \AttributeTok{row.names =} \ConstantTok{FALSE}\NormalTok{)}
  \FunctionTok{cat}\NormalTok{(}\StringTok{"Built and saved R\_hat\_df:"}\NormalTok{, }\FunctionTok{nrow}\NormalTok{(R\_hat\_df), }\StringTok{"rows x"}\NormalTok{,}
      \FunctionTok{ncol}\NormalTok{(R\_hat\_df), }\StringTok{"cols}\SpecialCharTok{\textbackslash{}n}\StringTok{"}\NormalTok{)}
\NormalTok{\}}
\end{Highlighting}
\end{Shaded}

\subsubsection{\texorpdfstring{\texttt{exposure\_matrix}: row-normalised
routing
shares}{exposure\_matrix: row-normalised routing shares}}\label{exposure_matrix-row-normalised-routing-shares}

\begin{Shaded}
\begin{Highlighting}[]
\ControlFlowTok{if}\NormalTok{ (}\FunctionTok{file.exists}\NormalTok{(}\FunctionTok{paste0}\NormalTok{(out\_dir, }\StringTok{"exposure\_matrix.csv"}\NormalTok{))) \{}

\NormalTok{  exposure\_matrix }\OtherTok{\textless{}{-}} \FunctionTok{read.csv}\NormalTok{(}\FunctionTok{paste0}\NormalTok{(out\_dir, }\StringTok{"exposure\_matrix.csv"}\NormalTok{),}
                               \AttributeTok{check.names =} \ConstantTok{FALSE}\NormalTok{,}
                               \AttributeTok{stringsAsFactors =} \ConstantTok{FALSE}\NormalTok{)}
  \FunctionTok{cat}\NormalTok{(}\StringTok{"Loaded exposure\_matrix from disk:"}\NormalTok{, }\FunctionTok{nrow}\NormalTok{(exposure\_matrix), }\StringTok{"rows}\SpecialCharTok{\textbackslash{}n}\StringTok{"}\NormalTok{)}

\NormalTok{\} }\ControlFlowTok{else}\NormalTok{ \{}

\NormalTok{  rowsums\_W       }\OtherTok{\textless{}{-}} \FunctionTok{rowSums}\NormalTok{(W\_raw)}
  \CommentTok{\# Avoid division by zero for any all{-}zero rows}
\NormalTok{  rowsums\_W[rowsums\_W }\SpecialCharTok{==} \DecValTok{0}\NormalTok{] }\OtherTok{\textless{}{-}} \DecValTok{1}
\NormalTok{  exposure\_raw    }\OtherTok{\textless{}{-}} \FunctionTok{sweep}\NormalTok{(W\_raw, }\DecValTok{1}\NormalTok{, rowsums\_W, }\StringTok{"/"}\NormalTok{)}

\NormalTok{  exposure\_matrix }\OtherTok{\textless{}{-}}\NormalTok{ r\_matrix\_idx }\SpecialCharTok{|\textgreater{}}
    \FunctionTok{select}\NormalTok{(mx\_state, mx\_municipality) }\SpecialCharTok{|\textgreater{}}
    \FunctionTok{bind\_cols}\NormalTok{(}\FunctionTok{as.data.frame}\NormalTok{(exposure\_raw))}

  \FunctionTok{write.csv}\NormalTok{(exposure\_matrix, }\FunctionTok{paste0}\NormalTok{(out\_dir, }\StringTok{"exposure\_matrix.csv"}\NormalTok{),}
            \AttributeTok{row.names =} \ConstantTok{FALSE}\NormalTok{)}
  \FunctionTok{cat}\NormalTok{(}\StringTok{"Built and saved exposure\_matrix:"}\NormalTok{, }\FunctionTok{nrow}\NormalTok{(exposure\_matrix), }\StringTok{"rows}\SpecialCharTok{\textbackslash{}n}\StringTok{"}\NormalTok{)}
\NormalTok{\}}
\end{Highlighting}
\end{Shaded}

\subsubsection{\texorpdfstring{\texttt{long\_ill\_df}: Illinois-only
long panel (all
municipalities)}{long\_ill\_df: Illinois-only long panel (all municipalities)}}\label{long_ill_df-illinois-only-long-panel-all-municipalities}

\begin{Shaded}
\begin{Highlighting}[]
\ControlFlowTok{if}\NormalTok{ (}\FunctionTok{file.exists}\NormalTok{(}\FunctionTok{paste0}\NormalTok{(out\_dir, }\StringTok{"long\_ill\_df.csv"}\NormalTok{))) \{}

\NormalTok{  long\_ill\_df }\OtherTok{\textless{}{-}} \FunctionTok{read.csv}\NormalTok{(}\FunctionTok{paste0}\NormalTok{(out\_dir, }\StringTok{"long\_ill\_df.csv"}\NormalTok{),}
                           \AttributeTok{stringsAsFactors =} \ConstantTok{FALSE}\NormalTok{) }\SpecialCharTok{|\textgreater{}}
    \FunctionTok{mutate}\NormalTok{(}\AttributeTok{year =} \FunctionTok{as.integer}\NormalTok{(year))}
  \FunctionTok{cat}\NormalTok{(}\StringTok{"Loaded long\_ill\_df from disk:"}\NormalTok{, }\FunctionTok{nrow}\NormalTok{(long\_ill\_df), }\StringTok{"rows}\SpecialCharTok{\textbackslash{}n}\StringTok{"}\NormalTok{)}

\NormalTok{\} }\ControlFlowTok{else}\NormalTok{ \{}

\NormalTok{  years       }\OtherTok{\textless{}{-}} \DecValTok{2010}\SpecialCharTok{:}\DecValTok{2024}
\NormalTok{  df\_list\_ill }\OtherTok{\textless{}{-}} \FunctionTok{list}\NormalTok{()}
  \ControlFlowTok{for}\NormalTok{ (yr }\ControlFlowTok{in}\NormalTok{ years) \{}
\NormalTok{    path  }\OtherTok{\textless{}{-}} \FunctionTok{file.path}\NormalTok{(us\_states\_path,}
                       \FunctionTok{paste0}\NormalTok{(}\StringTok{"Edos\_USA\_"}\NormalTok{, yr),}
                       \FunctionTok{paste0}\NormalTok{(}\StringTok{"ILLINOIS\_"}\NormalTok{, yr, }\StringTok{".xlsx"}\NormalTok{))}
\NormalTok{    dirty }\OtherTok{\textless{}{-}} \FunctionTok{read\_excel}\NormalTok{(path)}
\NormalTok{    clean }\OtherTok{\textless{}{-}} \FunctionTok{clean\_dfs}\NormalTok{(}\AttributeTok{df =}\NormalTok{ dirty, }\AttributeTok{state =} \StringTok{"Illinois"}\NormalTok{, }\AttributeTok{drop =} \ConstantTok{TRUE}\NormalTok{)}
\NormalTok{    df\_list\_ill[[}\FunctionTok{paste0}\NormalTok{(}\StringTok{"illinois\_"}\NormalTok{, yr)]] }\OtherTok{\textless{}{-}}\NormalTok{ clean}
\NormalTok{  \}}

  \CommentTok{\# Same deduplication fix: aggregate before joining}
\NormalTok{  wide\_ill }\OtherTok{\textless{}{-}} \FunctionTok{imap}\NormalTok{(df\_list\_ill, }\ControlFlowTok{function}\NormalTok{(df, key) \{}
\NormalTok{    df }\SpecialCharTok{|\textgreater{}}
      \FunctionTok{select}\NormalTok{(}\StringTok{\textasciigrave{}}\AttributeTok{origin state}\StringTok{\textasciigrave{}}\NormalTok{, }\StringTok{\textasciigrave{}}\AttributeTok{origin muni}\StringTok{\textasciigrave{}}\NormalTok{, weight) }\SpecialCharTok{|\textgreater{}}
      \FunctionTok{group\_by}\NormalTok{(}\StringTok{\textasciigrave{}}\AttributeTok{origin state}\StringTok{\textasciigrave{}}\NormalTok{, }\StringTok{\textasciigrave{}}\AttributeTok{origin muni}\StringTok{\textasciigrave{}}\NormalTok{) }\SpecialCharTok{|\textgreater{}}
      \FunctionTok{summarise}\NormalTok{(}\SpecialCharTok{!!}\AttributeTok{key :=} \FunctionTok{sum}\NormalTok{(weight, }\AttributeTok{na.rm =} \ConstantTok{TRUE}\NormalTok{), }\AttributeTok{.groups =} \StringTok{"drop"}\NormalTok{)}
\NormalTok{  \}) }\SpecialCharTok{|\textgreater{}}
    \FunctionTok{reduce}\NormalTok{(}\ControlFlowTok{function}\NormalTok{(a, b) }\FunctionTok{full\_join}\NormalTok{(a, b,}
                                    \AttributeTok{by =} \FunctionTok{c}\NormalTok{(}\StringTok{"origin state"}\NormalTok{, }\StringTok{"origin muni"}\NormalTok{))) }\SpecialCharTok{|\textgreater{}}
    \FunctionTok{mutate}\NormalTok{(}\FunctionTok{across}\NormalTok{(}\FunctionTok{where}\NormalTok{(is.numeric), }\SpecialCharTok{\textasciitilde{}}\FunctionTok{replace\_na}\NormalTok{(., }\DecValTok{0}\NormalTok{)))}

\NormalTok{  long\_ill\_df }\OtherTok{\textless{}{-}}\NormalTok{ wide\_ill }\SpecialCharTok{|\textgreater{}}
    \FunctionTok{pivot\_longer}\NormalTok{(}\AttributeTok{cols =} \SpecialCharTok{{-}}\FunctionTok{c}\NormalTok{(}\StringTok{\textasciigrave{}}\AttributeTok{origin state}\StringTok{\textasciigrave{}}\NormalTok{, }\StringTok{\textasciigrave{}}\AttributeTok{origin muni}\StringTok{\textasciigrave{}}\NormalTok{),}
                 \AttributeTok{names\_to =} \StringTok{"state\_year"}\NormalTok{, }\AttributeTok{values\_to =} \StringTok{"weight"}\NormalTok{) }\SpecialCharTok{|\textgreater{}}
    \FunctionTok{separate}\NormalTok{(state\_year, }\AttributeTok{into =} \FunctionTok{c}\NormalTok{(}\StringTok{"state"}\NormalTok{, }\StringTok{"year"}\NormalTok{),}
             \AttributeTok{sep =} \StringTok{"\_(?=}\SpecialCharTok{\textbackslash{}\textbackslash{}}\StringTok{d\{4\}$)"}\NormalTok{) }\SpecialCharTok{|\textgreater{}}
    \FunctionTok{mutate}\NormalTok{(}\AttributeTok{year =} \FunctionTok{as.integer}\NormalTok{(year))}

  \FunctionTok{write.csv}\NormalTok{(long\_ill\_df, }\FunctionTok{paste0}\NormalTok{(out\_dir, }\StringTok{"long\_ill\_df.csv"}\NormalTok{),}
            \AttributeTok{row.names =} \ConstantTok{FALSE}\NormalTok{)}
  \FunctionTok{cat}\NormalTok{(}\StringTok{"Built and saved long\_ill\_df:"}\NormalTok{, }\FunctionTok{nrow}\NormalTok{(long\_ill\_df), }\StringTok{"rows}\SpecialCharTok{\textbackslash{}n}\StringTok{"}\NormalTok{)}
\NormalTok{\}}
\end{Highlighting}
\end{Shaded}

\subsubsection{Derived views used in-paper (not saved
separately)}\label{derived-views-used-in-paper-not-saved-separately}

\begin{Shaded}
\begin{Highlighting}[]
\CommentTok{\# long\_joint\_df: long\_panel restricted to the 5 hand{-}picked bridge municipalities}
\NormalTok{joint\_pairs }\OtherTok{\textless{}{-}}\NormalTok{ tibble}\SpecialCharTok{::}\FunctionTok{tribble}\NormalTok{(}
  \SpecialCharTok{\textasciitilde{}}\StringTok{\textasciigrave{}}\AttributeTok{origin muni}\StringTok{\textasciigrave{}}\NormalTok{,        }\SpecialCharTok{\textasciitilde{}}\StringTok{\textasciigrave{}}\AttributeTok{origin state}\StringTok{\textasciigrave{}}\NormalTok{,}
  \StringTok{"Guadalajara"}\NormalTok{,         }\StringTok{"Jalisco"}\NormalTok{,}
  \StringTok{"Puebla"}\NormalTok{,              }\StringTok{"Puebla"}\NormalTok{,}
  \StringTok{"Atlixco"}\NormalTok{,             }\StringTok{"Puebla"}\NormalTok{,}
  \StringTok{"Acapulco De Juarez"}\NormalTok{,  }\StringTok{"Guerrero"}\NormalTok{,}
  \StringTok{"Gustavo A. Madero"}\NormalTok{,   }\StringTok{"Ciudad De Mexico"}
\NormalTok{)}

\NormalTok{long\_joint\_df }\OtherTok{\textless{}{-}}\NormalTok{ long\_panel }\SpecialCharTok{|\textgreater{}}
  \CommentTok{\# Rename to match the join keys used downstream}
  \FunctionTok{rename}\NormalTok{(}\StringTok{\textasciigrave{}}\AttributeTok{origin muni}\StringTok{\textasciigrave{}} \OtherTok{=} \StringTok{\textasciigrave{}}\AttributeTok{origin muni}\StringTok{\textasciigrave{}}\NormalTok{,}
         \StringTok{\textasciigrave{}}\AttributeTok{origin state}\StringTok{\textasciigrave{}} \OtherTok{=} \StringTok{\textasciigrave{}}\AttributeTok{origin state}\StringTok{\textasciigrave{}}\NormalTok{) }\SpecialCharTok{|\textgreater{}}
  \FunctionTok{inner\_join}\NormalTok{(joint\_pairs, }\AttributeTok{by =} \FunctionTok{c}\NormalTok{(}\StringTok{"origin muni"}\NormalTok{, }\StringTok{"origin state"}\NormalTok{))}
\end{Highlighting}
\end{Shaded}

\begin{center}\rule{0.5\linewidth}{0.5pt}\end{center}

\section{Introduction}\label{introduction}

Migration is not simply a bilateral relationship between a sending
community and a receiving country. It is a \textbf{network phenomenon}:
the same Mexican municipality often feeds multiple US states
simultaneously, and each US state draws from an overlapping set of
Mexican origin communities. This overlapping structure means that a
shock confined to one US state --- a recession, a policy change, a
fiscal crisis --- does not stay there. It travels backward through the
network to the Mexican municipalities exposed to that state, and then
forward along whatever other corridors those municipalities maintain,
eventually arriving in US states that were never directly touched by the
original shock.

This paper studies that transmission mechanism using the
\textbf{Illinois Budget Impasse} (July 1, 2015 -- July 10, 2017) as a
natural experiment. The impasse left Illinois without an operating
budget for 736 days --- the longest such gap in US history. It disrupted
state-funded construction, halted contractor payments, and froze social
services, hitting sectors where Mexican immigrants are concentrated.
Critically for identification, the impasse was:

\begin{enumerate}
\def\labelenumi{\arabic{enumi}.}
\tightlist
\item
  \textbf{Exogenous to migration decisions} --- it was a political
  stalemate over pension reform and tax policy, not an immigration
  policy, and Mexican migrants played no role in causing it;
\item
  \textbf{Geographically isolated} --- California and New York, the two
  other large-corridor states in our data, experienced no equivalent
  fiscal shock during this period; and
\item
  \textbf{Precisely timed} --- the start and end dates are known to the
  day, giving us a sharp pre/post window in annual MCAS data.
\end{enumerate}

Our contribution is to use the network's own structure --- bipartite
centrality, community membership, and the \(M^\top M\) state co-exposure
matrix --- to predict \emph{where} the shock should propagate before we
look at whether it did. This turns a purely empirical question into a
disciplined test: if the network predicts propagation from Illinois to
California through bridge municipalities, and if we then find a DiD
effect for California in exactly those municipalities with the highest
predicted exposure, the causal chain is credible.

\textbf{Data.} We use annual \emph{Matricula Consular} (MCAS) data from
2010--2024, which records the number of Matricula Consular IDs issued
per (US state, Mexican municipality) pair each year. These consular IDs
serve as a proxy for registered migrant presence. We complement this
with DS2 quarterly remittance outflow data per US state and DS1 actual
remittance inflow data at the Mexican municipality level from Banxico.
The migration weighting matrix \(W\) is a municipality \(\times\) state
matrix where entry \(W_{ij}\) is municipality \(i\)'s average share of
US state \(j\)'s total Matricula IDs, grand-total normalised across all
years (excluding COVID years 2020--2021) and averaged over the panel.

\textbf{Paper structure.} Section 2 constructs the bipartite migration
network, computes centrality measures that identify bridge
municipalities, and runs community detection on the state co-exposure
matrix to show that Illinois and California belong to the same migration
community. Section 3 defines the shock exposure variables and visualises
the impasse window in the time series. Section 4 presents the DiD
estimates: first the direct effect of Illinois exposure on imputed
remittances, then the second-order propagation test using California
flows through bridge municipalities. Section 5 discusses the cascade
interpretation, uses the TRUST Act as a positive shock robustness check,
and draws policy implications.

\begin{center}\rule{0.5\linewidth}{0.5pt}\end{center}

\section{Network Construction}\label{network-construction}

\subsection{The Bipartite Migration
Graph}\label{the-bipartite-migration-graph}

The migration network is naturally \textbf{bipartite}: nodes split into
two disjoint sets (Mexican municipalities \(\mathcal{M}\) and US states
\(\mathcal{S}\)), and edges only run between sets, never within them. An
edge \((i, j)\) with weight \(W_{ij}\) exists whenever municipality
\(i\) sends a positive share of migrants to state \(j\).

Formally, let \(W \in \mathbb{R}^{|\mathcal{M}| \times |\mathcal{S}|}\)
be the grand-total-normalised average migration weight matrix. The
bipartite graph \(G = (\mathcal{M} \cup \mathcal{S},\; E,\; W)\) has an
edge for each non-zero entry of \(W\), with that entry as the edge
weight.

\begin{Shaded}
\begin{Highlighting}[]
\CommentTok{\# Column{-}normalise W so each US state column sums to 1.}
\CommentTok{\# W\_col[i,j] = "of all migrants in state j, what share came from municipality i?"}
\CommentTok{\# This is column{-}stochastic and the right normalisation for the bipartite graph:}
\CommentTok{\# it weights each edge by how important municipality i is TO state j.}
\NormalTok{W\_col }\OtherTok{\textless{}{-}} \FunctionTok{sweep}\NormalTok{(W\_raw, }\DecValTok{2}\NormalTok{, }\FunctionTok{colSums}\NormalTok{(W\_raw), }\StringTok{"/"}\NormalTok{)}

\CommentTok{\# Build a long{-}format edge list from the weight matrix.}
\CommentTok{\# which(arr.ind=TRUE) returns row/col indices of all non{-}zero entries.}
\CommentTok{\# We threshold at a small epsilon to keep only edges with meaningful weight.}
\NormalTok{edge\_idx }\OtherTok{\textless{}{-}} \FunctionTok{which}\NormalTok{(W\_col }\SpecialCharTok{\textgreater{}} \FloatTok{1e{-}6}\NormalTok{, }\AttributeTok{arr.ind =} \ConstantTok{TRUE}\NormalTok{)}

\NormalTok{edge\_list }\OtherTok{\textless{}{-}} \FunctionTok{tibble}\NormalTok{(}
  \AttributeTok{from   =} \FunctionTok{rownames}\NormalTok{(W\_col)[edge\_idx[, }\DecValTok{1}\NormalTok{]],  }\CommentTok{\# Mexican municipality names}
  \AttributeTok{to     =} \FunctionTok{colnames}\NormalTok{(W\_col)[edge\_idx[, }\DecValTok{2}\NormalTok{]],  }\CommentTok{\# US state names}
  \AttributeTok{weight =}\NormalTok{ W\_col[edge\_idx]                  }\CommentTok{\# column{-}normalised weight}
\NormalTok{)}

\CommentTok{\# graph\_from\_data\_frame() builds the igraph object.}
\CommentTok{\# Columns 1 and 2 are edge endpoints; column 3 is the edge attribute.}
\NormalTok{G\_bip }\OtherTok{\textless{}{-}} \FunctionTok{graph\_from\_data\_frame}\NormalTok{(edge\_list, }\AttributeTok{directed =} \ConstantTok{FALSE}\NormalTok{)}

\CommentTok{\# Tag each vertex with its type: TRUE = Mexican municipality, FALSE = US state.}
\CommentTok{\# Assign directly from name membership — bipartite\_mapping() assigns by traversal}
\CommentTok{\# order, which is unpredictable and may flip the sets. rownames(W\_raw) are the}
\CommentTok{\# municipality names, so any vertex whose name appears there is a municipality.}
\FunctionTok{V}\NormalTok{(G\_bip)}\SpecialCharTok{$}\NormalTok{type }\OtherTok{\textless{}{-}} \FunctionTok{V}\NormalTok{(G\_bip)}\SpecialCharTok{$}\NormalTok{name }\SpecialCharTok{\%in\%} \FunctionTok{rownames}\NormalTok{(W\_raw)}

\FunctionTok{cat}\NormalTok{(}\StringTok{"Vertices:"}\NormalTok{, }\FunctionTok{vcount}\NormalTok{(G\_bip),}
    \StringTok{"("}\NormalTok{, }\FunctionTok{sum}\NormalTok{(}\SpecialCharTok{!}\FunctionTok{V}\NormalTok{(G\_bip)}\SpecialCharTok{$}\NormalTok{type), }\StringTok{"US states,"}\NormalTok{,}
    \FunctionTok{sum}\NormalTok{(}\FunctionTok{V}\NormalTok{(G\_bip)}\SpecialCharTok{$}\NormalTok{type), }\StringTok{"MX municipalities )}\SpecialCharTok{\textbackslash{}n}\StringTok{"}\NormalTok{)}
\FunctionTok{cat}\NormalTok{(}\StringTok{"Edges:"}\NormalTok{, }\FunctionTok{ecount}\NormalTok{(G\_bip), }\StringTok{"}\SpecialCharTok{\textbackslash{}n}\StringTok{"}\NormalTok{)}
\FunctionTok{cat}\NormalTok{(}\StringTok{"Is bipartite:"}\NormalTok{, }\FunctionTok{is\_bipartite}\NormalTok{(G\_bip), }\StringTok{"}\SpecialCharTok{\textbackslash{}n}\StringTok{"}\NormalTok{)}
\end{Highlighting}
\end{Shaded}

\begin{Shaded}
\begin{Highlighting}[]
\CommentTok{\# For visualisation, restrict to the three focal states and their top 15 municipalities.}
\NormalTok{focal\_states }\OtherTok{\textless{}{-}} \FunctionTok{c}\NormalTok{(}\StringTok{"Illinois"}\NormalTok{, }\StringTok{"California"}\NormalTok{, }\StringTok{"New York"}\NormalTok{)}
\NormalTok{top\_munis    }\OtherTok{\textless{}{-}}\NormalTok{ edge\_list }\SpecialCharTok{|\textgreater{}}
  \FunctionTok{filter}\NormalTok{(to }\SpecialCharTok{\%in\%}\NormalTok{ focal\_states) }\SpecialCharTok{|\textgreater{}}
  \FunctionTok{group\_by}\NormalTok{(from) }\SpecialCharTok{|\textgreater{}}
  \FunctionTok{summarise}\NormalTok{(}\AttributeTok{total\_weight =} \FunctionTok{sum}\NormalTok{(weight)) }\SpecialCharTok{|\textgreater{}}
  \FunctionTok{slice\_max}\NormalTok{(total\_weight, }\AttributeTok{n =} \DecValTok{15}\NormalTok{) }\SpecialCharTok{|\textgreater{}}
  \FunctionTok{pull}\NormalTok{(from)}

\CommentTok{\# Subset the edge list to the focal states and top municipalities}
\NormalTok{edge\_sub }\OtherTok{\textless{}{-}}\NormalTok{ edge\_list }\SpecialCharTok{|\textgreater{}}
  \FunctionTok{filter}\NormalTok{(from }\SpecialCharTok{\%in\%}\NormalTok{ top\_munis, to }\SpecialCharTok{\%in\%}\NormalTok{ focal\_states)}

\NormalTok{G\_sub }\OtherTok{\textless{}{-}} \FunctionTok{graph\_from\_data\_frame}\NormalTok{(edge\_sub, }\AttributeTok{directed =} \ConstantTok{FALSE}\NormalTok{)}
\FunctionTok{V}\NormalTok{(G\_sub)}\SpecialCharTok{$}\NormalTok{type }\OtherTok{\textless{}{-}} \FunctionTok{V}\NormalTok{(G\_sub)}\SpecialCharTok{$}\NormalTok{name }\SpecialCharTok{\%in\%}\NormalTok{ top\_munis  }\CommentTok{\# TRUE = municipality}

\CommentTok{\# Color nodes by type: municipalities in teal, states in orange}
\FunctionTok{V}\NormalTok{(G\_sub)}\SpecialCharTok{$}\NormalTok{color }\OtherTok{\textless{}{-}} \FunctionTok{ifelse}\NormalTok{(}\FunctionTok{V}\NormalTok{(G\_sub)}\SpecialCharTok{$}\NormalTok{type, }\StringTok{"\#4DADA7"}\NormalTok{, }\StringTok{"\#E07B39"}\NormalTok{)}
\FunctionTok{V}\NormalTok{(G\_sub)}\SpecialCharTok{$}\NormalTok{size  }\OtherTok{\textless{}{-}} \FunctionTok{ifelse}\NormalTok{(}\FunctionTok{V}\NormalTok{(G\_sub)}\SpecialCharTok{$}\NormalTok{type, }\DecValTok{8}\NormalTok{, }\DecValTok{14}\NormalTok{)}

\CommentTok{\# Scale edge width by weight}
\FunctionTok{E}\NormalTok{(G\_sub)}\SpecialCharTok{$}\NormalTok{width }\OtherTok{\textless{}{-}} \FunctionTok{E}\NormalTok{(G\_sub)}\SpecialCharTok{$}\NormalTok{weight }\SpecialCharTok{*} \DecValTok{80}

\FunctionTok{plot}\NormalTok{(}
\NormalTok{  G\_sub,}
  \AttributeTok{layout            =} \FunctionTok{layout\_as\_bipartite}\NormalTok{(G\_sub),}
  \AttributeTok{vertex.label      =} \FunctionTok{V}\NormalTok{(G\_sub)}\SpecialCharTok{$}\NormalTok{name,}
  \AttributeTok{vertex.label.cex  =} \FloatTok{0.55}\NormalTok{,}
  \AttributeTok{vertex.label.dist =} \FloatTok{0.8}\NormalTok{,}
  \AttributeTok{vertex.color      =} \FunctionTok{V}\NormalTok{(G\_sub)}\SpecialCharTok{$}\NormalTok{color,}
  \AttributeTok{vertex.size       =} \FunctionTok{V}\NormalTok{(G\_sub)}\SpecialCharTok{$}\NormalTok{size,}
  \AttributeTok{edge.width        =} \FunctionTok{E}\NormalTok{(G\_sub)}\SpecialCharTok{$}\NormalTok{width,}
  \AttributeTok{edge.color        =} \FunctionTok{adjustcolor}\NormalTok{(}\StringTok{"grey40"}\NormalTok{, }\AttributeTok{alpha.f =} \FloatTok{0.5}\NormalTok{),}
  \AttributeTok{main =} \StringTok{"Migration Network: Top 15 MX Municipalities to IL, CA, NY"}
\NormalTok{)}
\FunctionTok{legend}\NormalTok{(}\StringTok{"bottomright"}\NormalTok{,}
       \AttributeTok{legend =} \FunctionTok{c}\NormalTok{(}\StringTok{"MX Municipality"}\NormalTok{, }\StringTok{"US State"}\NormalTok{),}
       \AttributeTok{fill   =} \FunctionTok{c}\NormalTok{(}\StringTok{"\#4DADA7"}\NormalTok{, }\StringTok{"\#E07B39"}\NormalTok{),}
       \AttributeTok{bty    =} \StringTok{"n"}\NormalTok{, }\AttributeTok{cex =} \FloatTok{0.8}\NormalTok{)}
\end{Highlighting}
\end{Shaded}

\subsection{Centrality Analysis}\label{centrality-analysis}

Centrality measures tell us which municipalities are structurally
important in the network --- not just because they send many migrants to
one state, but because of their position relative to the whole graph.
Two measures are especially relevant here.

\textbf{Betweenness centrality} counts how often a node lies on the
shortest path between two other nodes. In our bipartite graph, a
municipality with high betweenness centrality frequently lies on the
shortest path between two US states, meaning it is the natural relay
node between those states. These are exactly the bridge municipalities
through which an Illinois shock would reach California: they sit
\emph{between} the two states in network space.

\textbf{Eigenvector centrality} scores a node by the centrality of its
neighbors. A municipality that sends migrants to several large,
well-connected states gets a high eigenvector centrality score, even if
its raw weight to any single state is modest. This makes it a better
continuous measure of treatment intensity than a simple binary high/low
IL-exposure cutoff.

\begin{Shaded}
\begin{Highlighting}[]
\CommentTok{\# Betweenness centrality: how often does node i sit on shortest paths between other nodes?}
\CommentTok{\# normalized=TRUE divides by the maximum possible betweenness so scores are in [0,1].}
\CommentTok{\# weights=NA tells igraph to ignore edge weights for path{-}length computation}
\CommentTok{\# (otherwise higher{-}weight edges are treated as SHORTER, which is semantically wrong here).}
\NormalTok{betweenness\_scores }\OtherTok{\textless{}{-}} \FunctionTok{betweenness}\NormalTok{(G\_bip, }\AttributeTok{normalized =} \ConstantTok{TRUE}\NormalTok{, }\AttributeTok{weights =} \ConstantTok{NA}\NormalTok{)}

\CommentTok{\# Eigenvector centrality: scored by the centrality of a node\textquotesingle{}s neighbors.}
\CommentTok{\# $vector returns the centrality scores (the leading eigenvector of the adjacency matrix).}
\NormalTok{eigen\_scores }\OtherTok{\textless{}{-}} \FunctionTok{eigen\_centrality}\NormalTok{(G\_bip, }\AttributeTok{weights =} \FunctionTok{E}\NormalTok{(G\_bip)}\SpecialCharTok{$}\NormalTok{weight)}\SpecialCharTok{$}\NormalTok{vector}

\CommentTok{\# Collect into a tidy data frame, restricting to MX municipalities only}
\NormalTok{centrality\_df }\OtherTok{\textless{}{-}} \FunctionTok{tibble}\NormalTok{(}
  \AttributeTok{name              =} \FunctionTok{V}\NormalTok{(G\_bip)}\SpecialCharTok{$}\NormalTok{name,}
  \AttributeTok{is\_municipality   =} \FunctionTok{V}\NormalTok{(G\_bip)}\SpecialCharTok{$}\NormalTok{type,}
  \AttributeTok{betweenness       =}\NormalTok{ betweenness\_scores,}
  \AttributeTok{eigenvector       =}\NormalTok{ eigen\_scores}
\NormalTok{) }\SpecialCharTok{|\textgreater{}}
  \FunctionTok{filter}\NormalTok{(is\_municipality) }\SpecialCharTok{|\textgreater{}}   \CommentTok{\# keep only MX municipality nodes}
  \FunctionTok{arrange}\NormalTok{(}\FunctionTok{desc}\NormalTok{(betweenness))}

\FunctionTok{head}\NormalTok{(centrality\_df, }\DecValTok{20}\NormalTok{)   }\CommentTok{\# top 20 municipalities by betweenness}
\end{Highlighting}
\end{Shaded}

\begin{Shaded}
\begin{Highlighting}[]
\CommentTok{\# Municipalities with the highest betweenness are the most likely transmission nodes}
\CommentTok{\# for cross{-}state shock propagation. We label these as our "bridge municipalities."}
\NormalTok{top20\_between }\OtherTok{\textless{}{-}}\NormalTok{ centrality\_df }\SpecialCharTok{|\textgreater{}} \FunctionTok{slice\_max}\NormalTok{(betweenness, }\AttributeTok{n =} \DecValTok{20}\NormalTok{)}

\FunctionTok{ggplot}\NormalTok{(top20\_between,}
       \FunctionTok{aes}\NormalTok{(}\AttributeTok{x =}\NormalTok{ betweenness, }\AttributeTok{y =} \FunctionTok{reorder}\NormalTok{(name, betweenness))) }\SpecialCharTok{+}
  \FunctionTok{geom\_col}\NormalTok{(}\AttributeTok{fill =} \StringTok{"\#4DADA7"}\NormalTok{) }\SpecialCharTok{+}
  \FunctionTok{labs}\NormalTok{(}\AttributeTok{x =} \StringTok{"Betweenness Centrality (normalised)"}\NormalTok{,}
       \AttributeTok{y =} \ConstantTok{NULL}\NormalTok{,}
       \AttributeTok{title =} \StringTok{"Top 20 MX Municipalities by Betweenness Centrality"}\NormalTok{,}
       \AttributeTok{subtitle =} \StringTok{"High betweenness = bridge node between US states"}\NormalTok{) }\SpecialCharTok{+}
  \FunctionTok{theme\_minimal}\NormalTok{(}\AttributeTok{base\_size =} \DecValTok{9}\NormalTok{)}
\end{Highlighting}
\end{Shaded}

\begin{Shaded}
\begin{Highlighting}[]
\CommentTok{\# Scatter plot of the two centrality measures.}
\CommentTok{\# Top{-}right: both central AND well{-}connected {-}{-} the most important bridge nodes.}
\CommentTok{\# Top{-}left:  high betweenness but low eigenvector {-}{-} bridges to peripheral states.}
\CommentTok{\# Bottom{-}right: well{-}connected but not a bridge {-}{-} large senders to one dominant state.}
\FunctionTok{ggplot}\NormalTok{(centrality\_df, }\FunctionTok{aes}\NormalTok{(}\AttributeTok{x =}\NormalTok{ eigenvector, }\AttributeTok{y =}\NormalTok{ betweenness)) }\SpecialCharTok{+}
  \FunctionTok{geom\_point}\NormalTok{(}\AttributeTok{alpha =} \FloatTok{0.35}\NormalTok{, }\AttributeTok{size =} \FloatTok{0.9}\NormalTok{, }\AttributeTok{color =} \StringTok{"steelblue"}\NormalTok{) }\SpecialCharTok{+}
  \CommentTok{\# Label the top 10 by betweenness}
\NormalTok{  ggrepel}\SpecialCharTok{::}\FunctionTok{geom\_text\_repel}\NormalTok{(}
    \AttributeTok{data    =} \FunctionTok{slice\_max}\NormalTok{(centrality\_df, betweenness, }\AttributeTok{n =} \DecValTok{10}\NormalTok{),}
    \FunctionTok{aes}\NormalTok{(}\AttributeTok{label =}\NormalTok{ name),}
    \AttributeTok{size    =} \FloatTok{2.5}\NormalTok{,}
    \AttributeTok{max.overlaps =} \DecValTok{15}
\NormalTok{  ) }\SpecialCharTok{+}
  \FunctionTok{labs}\NormalTok{(}\AttributeTok{x =} \StringTok{"Eigenvector Centrality"}\NormalTok{,}
       \AttributeTok{y =} \StringTok{"Betweenness Centrality"}\NormalTok{,}
       \AttributeTok{title =} \StringTok{"Centrality in the MX{-}US Migration Network"}\NormalTok{) }\SpecialCharTok{+}
  \FunctionTok{theme\_minimal}\NormalTok{(}\AttributeTok{base\_size =} \DecValTok{9}\NormalTok{)}
\end{Highlighting}
\end{Shaded}

The top-betweenness municipalities are our operationalisation of
``bridge nodes.'' We formally define a binary bridge indicator and a
continuous treatment intensity variable:

\begin{Shaded}
\begin{Highlighting}[]
\CommentTok{\# Define bridge municipalities as those above the 90th percentile of betweenness.}
\CommentTok{\# This threshold is deliberately conservative: we want nodes that are unambiguously}
\CommentTok{\# structural bridges, not marginal cases.}
\NormalTok{p90\_between }\OtherTok{\textless{}{-}} \FunctionTok{quantile}\NormalTok{(centrality\_df}\SpecialCharTok{$}\NormalTok{betweenness, }\FloatTok{0.90}\NormalTok{)}

\NormalTok{centrality\_df }\OtherTok{\textless{}{-}}\NormalTok{ centrality\_df }\SpecialCharTok{|\textgreater{}}
  \FunctionTok{mutate}\NormalTok{(}
    \AttributeTok{is\_bridge         =}\NormalTok{ betweenness }\SpecialCharTok{\textgreater{}}\NormalTok{ p90\_between,   }\CommentTok{\# binary bridge indicator}
    \CommentTok{\# Eigenvector centrality as continuous treatment intensity for the DiD}
    \CommentTok{\# (standardised to mean 0, SD 1 for interpretable coefficients)}
    \AttributeTok{eigen\_std         =} \FunctionTok{scale}\NormalTok{(eigenvector)[, }\DecValTok{1}\NormalTok{]}
\NormalTok{  )}

\CommentTok{\# Summary: how many bridge municipalities? What states do they connect?}
\FunctionTok{cat}\NormalTok{(}\StringTok{"Bridge municipalities (top{-}10\% betweenness):"}\NormalTok{, }\FunctionTok{sum}\NormalTok{(centrality\_df}\SpecialCharTok{$}\NormalTok{is\_bridge), }\StringTok{"}\SpecialCharTok{\textbackslash{}n}\StringTok{"}\NormalTok{)}
\FunctionTok{cat}\NormalTok{(}\StringTok{"90th percentile betweenness threshold:"}\NormalTok{, }\FunctionTok{round}\NormalTok{(p90\_between, }\DecValTok{5}\NormalTok{), }\StringTok{"}\SpecialCharTok{\textbackslash{}n}\StringTok{"}\NormalTok{)}
\end{Highlighting}
\end{Shaded}

\subsection{\texorpdfstring{Community Detection on the \(M^\top M\)
State Co-Exposure
Matrix}{Community Detection on the M\^{}\textbackslash top M State Co-Exposure Matrix}}\label{community-detection-on-the-mtop-m-state-co-exposure-matrix}

A municipality sends migrants to multiple US states. When two US states
draw heavily from the \emph{same} set of municipalities, a shock in one
state will be transmitted to the other via those shared sending
communities. We can measure this shared-origin overlap directly: compute
\(M^\top M\) (the inner product of the weight matrix with itself), which
gives a \textbf{state \(\times\) state} symmetric matrix where entry
\((j, k)\) is proportional to the cosine similarity between state
\(j\)'s and state \(k\)'s municipality distributions.

\[[M^\top M]_{jk} = \sum_{i \in \mathcal{M}} W_{ij} \cdot W_{ik}\]

A large \((j,k)\) entry means states \(j\) and \(k\) recruit from many
of the same Mexican municipalities --- they are in direct competition
for the same sending communities and are maximally exposed to shared
shocks. We treat \(M^\top M\) as a weighted adjacency matrix between US
states and run \textbf{Louvain community detection} on it. The
communities it finds are groups of US states that compete for the same
Mexican sending municipalities. If Illinois and California belong to the
same community, indirect shock propagation is structurally plausible;
states in \emph{different} communities from Illinois should show no
propagation.

\begin{Shaded}
\begin{Highlighting}[]
\CommentTok{\# M = W\_col (column{-}normalised weight matrix, n\_munis x n\_states)}
\CommentTok{\# M\textquotesingle{}M = t(W\_col) \%*\% W\_col {-}{-} (n\_states x n\_states) co{-}exposure matrix}
\CommentTok{\# Each entry (j,k) = sum over munis of W[i,j] * W[i,k]}
\CommentTok{\# = dot product of state j\textquotesingle{}s and state k\textquotesingle{}s municipality weight vectors}
\CommentTok{\# = how much overlap the two states have in their sending communities}

\NormalTok{MtM }\OtherTok{\textless{}{-}} \FunctionTok{t}\NormalTok{(W\_col) }\SpecialCharTok{\%*\%}\NormalTok{ W\_col   }\CommentTok{\# matrix multiply: (n\_states x n\_munis) \%*\% (n\_munis x n\_states)}

\CommentTok{\# Symmetry check: M\textquotesingle{}M should be exactly symmetric}
\FunctionTok{cat}\NormalTok{(}\StringTok{"Max asymmetry:"}\NormalTok{, }\FunctionTok{max}\NormalTok{(}\FunctionTok{abs}\NormalTok{(MtM }\SpecialCharTok{{-}} \FunctionTok{t}\NormalTok{(MtM))), }\StringTok{"}\SpecialCharTok{\textbackslash{}n}\StringTok{"}\NormalTok{)  }\CommentTok{\# should be \textasciitilde{}0}

\CommentTok{\# Zero out the diagonal (self{-}overlap is not meaningful for community detection)}
\FunctionTok{diag}\NormalTok{(MtM) }\OtherTok{\textless{}{-}} \DecValTok{0}

\CommentTok{\# Inspect the top co{-}exposed state pairs}
\NormalTok{MtM\_df }\OtherTok{\textless{}{-}} \FunctionTok{as.data.frame}\NormalTok{(}\FunctionTok{as.table}\NormalTok{(MtM)) }\SpecialCharTok{|\textgreater{}}
  \FunctionTok{rename}\NormalTok{(}\AttributeTok{state\_a =}\NormalTok{ Var1, }\AttributeTok{state\_b =}\NormalTok{ Var2, }\AttributeTok{overlap =}\NormalTok{ Freq) }\SpecialCharTok{|\textgreater{}}
  \FunctionTok{filter}\NormalTok{(state\_a }\SpecialCharTok{\textless{}}\NormalTok{ state\_b) }\SpecialCharTok{|\textgreater{}}    \CommentTok{\# keep only upper triangle (symmetric matrix)}
  \FunctionTok{arrange}\NormalTok{(}\FunctionTok{desc}\NormalTok{(overlap))}

\FunctionTok{head}\NormalTok{(MtM\_df, }\DecValTok{15}\NormalTok{)   }\CommentTok{\# which state pairs share the most sending communities?}
\end{Highlighting}
\end{Shaded}

\begin{Shaded}
\begin{Highlighting}[]
\CommentTok{\# Build an igraph graph from M\textquotesingle{}M treated as a weighted adjacency matrix.}
\CommentTok{\# graph\_from\_adjacency\_matrix() reads a square matrix as an adjacency matrix.}
\CommentTok{\# mode="undirected" because M\textquotesingle{}M is symmetric.}
\CommentTok{\# weighted=TRUE uses the matrix values as edge weights.}
\NormalTok{G\_states }\OtherTok{\textless{}{-}} \FunctionTok{graph\_from\_adjacency\_matrix}\NormalTok{(}
\NormalTok{  MtM,}
  \AttributeTok{mode     =} \StringTok{"undirected"}\NormalTok{,}
  \AttributeTok{weighted =} \ConstantTok{TRUE}\NormalTok{,}
  \AttributeTok{diag     =} \ConstantTok{FALSE}          \CommentTok{\# no self{-}loops}
\NormalTok{)}

\CommentTok{\# Run Louvain community detection.}
\CommentTok{\# Louvain is a greedy modularity{-}maximisation algorithm: it repeatedly merges nodes}
\CommentTok{\# into communities to maximise the fraction of within{-}community edge weight}
\CommentTok{\# relative to what would be expected under a null model.}
\CommentTok{\# set.seed() ensures reproducibility (Louvain has a stochastic component).}
\FunctionTok{set.seed}\NormalTok{(}\DecValTok{42}\NormalTok{)}
\NormalTok{communities }\OtherTok{\textless{}{-}} \FunctionTok{cluster\_louvain}\NormalTok{(G\_states, }\AttributeTok{weights =} \FunctionTok{E}\NormalTok{(G\_states)}\SpecialCharTok{$}\NormalTok{weight)}

\CommentTok{\# Print community assignments}
\NormalTok{community\_df }\OtherTok{\textless{}{-}} \FunctionTok{tibble}\NormalTok{(}
  \AttributeTok{state     =} \FunctionTok{V}\NormalTok{(G\_states)}\SpecialCharTok{$}\NormalTok{name,}
  \AttributeTok{community =} \FunctionTok{membership}\NormalTok{(communities)   }\CommentTok{\# integer community label for each state}
\NormalTok{) }\SpecialCharTok{|\textgreater{}}
  \FunctionTok{arrange}\NormalTok{(community, state)}

\FunctionTok{cat}\NormalTok{(}\StringTok{"Number of communities detected:"}\NormalTok{, }\FunctionTok{max}\NormalTok{(community\_df}\SpecialCharTok{$}\NormalTok{community), }\StringTok{"}\SpecialCharTok{\textbackslash{}n}\StringTok{"}\NormalTok{)}
\FunctionTok{cat}\NormalTok{(}\StringTok{"Modularity:"}\NormalTok{, }\FunctionTok{round}\NormalTok{(}\FunctionTok{modularity}\NormalTok{(communities), }\DecValTok{4}\NormalTok{), }\StringTok{"}\SpecialCharTok{\textbackslash{}n\textbackslash{}n}\StringTok{"}\NormalTok{)}

\CommentTok{\# Check which community Illinois and California belong to}
\NormalTok{community\_df }\SpecialCharTok{|\textgreater{}} \FunctionTok{filter}\NormalTok{(state }\SpecialCharTok{\%in\%} \FunctionTok{c}\NormalTok{(}\StringTok{"Illinois"}\NormalTok{, }\StringTok{"California"}\NormalTok{, }\StringTok{"New York"}\NormalTok{))}
\end{Highlighting}
\end{Shaded}

\begin{Shaded}
\begin{Highlighting}[]
\CommentTok{\# Assign community colors to nodes}
\NormalTok{n\_comm   }\OtherTok{\textless{}{-}} \FunctionTok{max}\NormalTok{(community\_df}\SpecialCharTok{$}\NormalTok{community)}
\NormalTok{pal      }\OtherTok{\textless{}{-}}\NormalTok{ RColorBrewer}\SpecialCharTok{::}\FunctionTok{brewer.pal}\NormalTok{(}\FunctionTok{min}\NormalTok{(n\_comm, }\DecValTok{9}\NormalTok{), }\StringTok{"Set1"}\NormalTok{)}
\NormalTok{node\_col }\OtherTok{\textless{}{-}}\NormalTok{ pal[}\FunctionTok{membership}\NormalTok{(communities)]}

\CommentTok{\# Threshold edges so the plot is readable: only show the top 10\% strongest overlaps}
\NormalTok{q90\_edge }\OtherTok{\textless{}{-}} \FunctionTok{quantile}\NormalTok{(}\FunctionTok{E}\NormalTok{(G\_states)}\SpecialCharTok{$}\NormalTok{weight, }\FloatTok{0.90}\NormalTok{)}
\NormalTok{G\_plot   }\OtherTok{\textless{}{-}} \FunctionTok{delete\_edges}\NormalTok{(G\_states, }\FunctionTok{which}\NormalTok{(}\FunctionTok{E}\NormalTok{(G\_states)}\SpecialCharTok{$}\NormalTok{weight }\SpecialCharTok{\textless{}}\NormalTok{ q90\_edge))}

\CommentTok{\# Mark the three focal states with a different shape}
\FunctionTok{V}\NormalTok{(G\_plot)}\SpecialCharTok{$}\NormalTok{shape }\OtherTok{\textless{}{-}} \FunctionTok{ifelse}\NormalTok{(}
  \FunctionTok{V}\NormalTok{(G\_plot)}\SpecialCharTok{$}\NormalTok{name }\SpecialCharTok{\%in\%} \FunctionTok{c}\NormalTok{(}\StringTok{"Illinois"}\NormalTok{, }\StringTok{"California"}\NormalTok{, }\StringTok{"New York"}\NormalTok{),}
  \StringTok{"square"}\NormalTok{, }\StringTok{"circle"}
\NormalTok{)}

\FunctionTok{plot}\NormalTok{(}
\NormalTok{  G\_plot,}
  \AttributeTok{layout         =} \FunctionTok{layout\_with\_fr}\NormalTok{(G\_plot),   }\CommentTok{\# Fruchterman{-}Reingold force{-}directed layout}
  \AttributeTok{vertex.color   =}\NormalTok{ node\_col[}\FunctionTok{match}\NormalTok{(}\FunctionTok{V}\NormalTok{(G\_plot)}\SpecialCharTok{$}\NormalTok{name, }\FunctionTok{V}\NormalTok{(G\_states)}\SpecialCharTok{$}\NormalTok{name)],}
  \AttributeTok{vertex.size    =} \DecValTok{7}\NormalTok{,}
  \AttributeTok{vertex.shape   =} \FunctionTok{V}\NormalTok{(G\_plot)}\SpecialCharTok{$}\NormalTok{shape,}
  \AttributeTok{vertex.label   =} \FunctionTok{V}\NormalTok{(G\_plot)}\SpecialCharTok{$}\NormalTok{name,}
  \AttributeTok{vertex.label.cex =} \FloatTok{0.4}\NormalTok{,}
  \AttributeTok{edge.width     =} \FunctionTok{E}\NormalTok{(G\_plot)}\SpecialCharTok{$}\NormalTok{weight }\SpecialCharTok{*} \DecValTok{40}\NormalTok{,}
  \AttributeTok{edge.color     =} \FunctionTok{adjustcolor}\NormalTok{(}\StringTok{"grey60"}\NormalTok{, }\AttributeTok{alpha.f =} \FloatTok{0.4}\NormalTok{),}
  \AttributeTok{main =} \StringTok{"State Co{-}Exposure Network: Louvain Communities}\SpecialCharTok{\textbackslash{}n}\StringTok{(squares = IL, CA, NY)"}
\NormalTok{)}
\end{Highlighting}
\end{Shaded}

\begin{Shaded}
\begin{Highlighting}[]
\CommentTok{\# Print full community membership table}
\NormalTok{community\_df }\SpecialCharTok{|\textgreater{}}
  \FunctionTok{group\_by}\NormalTok{(community) }\SpecialCharTok{|\textgreater{}}
  \FunctionTok{summarise}\NormalTok{(}
    \AttributeTok{n\_states =} \FunctionTok{n}\NormalTok{(),}
    \AttributeTok{states   =} \FunctionTok{paste}\NormalTok{(}\FunctionTok{sort}\NormalTok{(state), }\AttributeTok{collapse =} \StringTok{", "}\NormalTok{)}
\NormalTok{  ) }\SpecialCharTok{|\textgreater{}}
\NormalTok{  knitr}\SpecialCharTok{::}\FunctionTok{kable}\NormalTok{(}
    \AttributeTok{caption =} \StringTok{"Louvain community membership: US states grouped by shared MX sending communities."}\NormalTok{,}
    \AttributeTok{booktabs =} \ConstantTok{TRUE}
\NormalTok{  )}
\end{Highlighting}
\end{Shaded}

The community detection result provides the key structural motivation
for our propagation hypothesis. If Illinois and California are assigned
to the same community --- meaning they compete for the same Mexican
sending municipalities --- then any disruption to Illinois-bound
migration will necessarily reach municipalities that also feed
California, making cross-state propagation structurally inevitable.
States in \emph{different} communities serve as a natural placebo group:
they should show no propagation, even if they experience their own
shocks, because they draw from a disjoint pool of sending communities.

\begin{center}\rule{0.5\linewidth}{0.5pt}\end{center}

\section{Shock Identification}\label{shock-identification}

\subsection{The Illinois Budget Impasse as a Natural
Experiment}\label{the-illinois-budget-impasse-as-a-natural-experiment}

The Illinois Budget Impasse ran from July 1, 2015 through July 10, 2017
--- 736 days. Without an operating budget, the state could not fund
construction projects, pay contractors, or maintain social service
contracts. These are precisely the labor markets where low-wage
immigrant workers, including Mexican nationals, are concentrated.

For causal identification, three features make this event attractive:

\textbf{Exogeneity.} The impasse originated in a structural conflict
between Governor Rauner (who demanded union-weakening legislation and a
property tax freeze as conditions for any budget) and a Democratic
supermajority legislature that refused those terms. The deadlock was
entirely about domestic fiscal and labor politics. Mexican migration
decisions played no role in causing the crisis and were not targeted by
any policy within it.

\textbf{Specificity.} California and New York --- the other
large-recipient states in our data --- had no equivalent fiscal event
during 2015--2017. California was in a period of fiscal expansion, and
New York showed no major budget disruptions. This makes a
parallel-trends assumption between Illinois and the control states
substantially more credible than a typical DiD design where treatment
and control states may have correlated unobservables.

\textbf{Sharp timing.} The impasse began on a specific date (July 1,
2015, when the new fiscal year started without a passed budget) and
ended on a specific date (July 10, 2017, when the legislature overrode
Rauner's veto). In annual MCAS data, this maps cleanly to treatment
years 2015 and 2016, with 2017 as a partial-treatment year and 2018
onward as the post-recovery window.

\subsection{The Exposure Matrix}\label{the-exposure-matrix}

Not all municipalities are equally affected by an Illinois shock. The
key variable is each municipality's \textbf{Illinois routing
probability} from the row-normalised exposure matrix: of all the
migrants from municipality \(i\), what fraction historically went to
Illinois?

\[\text{Exposure}_i^{IL} = \frac{W_{i,IL}}{\sum_j W_{ij}} = \text{ExposureMatrix}[i, \text{Illinois}]\]

This gives a continuous measure of pre-shock Illinois dependence. We
complement it with the eigenvector centrality score computed in Section
2, which captures not just dependence on Illinois but the municipality's
overall structural importance in the network.

\begin{Shaded}
\begin{Highlighting}[]
\CommentTok{\# Row{-}normalised exposure matrix (each row sums to 1)}
\CommentTok{\# exposure\_matrix already computed in the data pipeline}
\NormalTok{rowsums\_W     }\OtherTok{\textless{}{-}} \FunctionTok{rowSums}\NormalTok{(W\_raw)}
\NormalTok{exposure\_raw  }\OtherTok{\textless{}{-}} \FunctionTok{sweep}\NormalTok{(W\_raw, }\DecValTok{1}\NormalTok{, rowsums\_W, }\StringTok{"/"}\NormalTok{)   }\CommentTok{\# row{-}normalise}

\NormalTok{exposure\_matrix }\OtherTok{\textless{}{-}}\NormalTok{ r\_matrix\_idx }\SpecialCharTok{|\textgreater{}}
  \FunctionTok{select}\NormalTok{(mx\_state, mx\_municipality) }\SpecialCharTok{|\textgreater{}}
  \FunctionTok{bind\_cols}\NormalTok{(}\FunctionTok{as.data.frame}\NormalTok{(exposure\_raw))}

\CommentTok{\# Extract Illinois{-}specific exposure for each municipality}
\NormalTok{il\_exposure }\OtherTok{\textless{}{-}}\NormalTok{ exposure\_matrix }\SpecialCharTok{|\textgreater{}}
  \FunctionTok{select}\NormalTok{(mx\_state, mx\_municipality, }\AttributeTok{il\_exposure =}\NormalTok{ Illinois) }\SpecialCharTok{|\textgreater{}}
  \CommentTok{\# Join centrality scores}
  \FunctionTok{left\_join}\NormalTok{(}
\NormalTok{    centrality\_df }\SpecialCharTok{|\textgreater{}} \FunctionTok{select}\NormalTok{(name, betweenness, eigenvector, eigen\_std, is\_bridge),}
    \AttributeTok{by =} \FunctionTok{c}\NormalTok{(}\StringTok{"mx\_municipality"} \OtherTok{=} \StringTok{"name"}\NormalTok{)}
\NormalTok{  )}

\CommentTok{\# Distribution of IL exposure}
\FunctionTok{cat}\NormalTok{(}\StringTok{"IL exposure quantiles:}\SpecialCharTok{\textbackslash{}n}\StringTok{"}\NormalTok{)}
\FunctionTok{print}\NormalTok{(}\FunctionTok{quantile}\NormalTok{(il\_exposure}\SpecialCharTok{$}\NormalTok{il\_exposure, }\FunctionTok{c}\NormalTok{(}\FloatTok{0.25}\NormalTok{, }\FloatTok{0.50}\NormalTok{, }\FloatTok{0.75}\NormalTok{, }\FloatTok{0.90}\NormalTok{, }\FloatTok{0.95}\NormalTok{, }\FloatTok{0.99}\NormalTok{),}
               \AttributeTok{na.rm =} \ConstantTok{TRUE}\NormalTok{))}

\CommentTok{\# Define treatment group: municipalities with above{-}median IL exposure}
\CommentTok{\# We use median split as the primary binary treatment definition}
\CommentTok{\# and eigenvector centrality as a continuous alternative}
\NormalTok{med\_il }\OtherTok{\textless{}{-}} \FunctionTok{median}\NormalTok{(il\_exposure}\SpecialCharTok{$}\NormalTok{il\_exposure[il\_exposure}\SpecialCharTok{$}\NormalTok{il\_exposure }\SpecialCharTok{\textgreater{}} \DecValTok{0}\NormalTok{], }\AttributeTok{na.rm =} \ConstantTok{TRUE}\NormalTok{)}
\NormalTok{il\_exposure }\OtherTok{\textless{}{-}}\NormalTok{ il\_exposure }\SpecialCharTok{|\textgreater{}}
  \FunctionTok{mutate}\NormalTok{(}
    \AttributeTok{treated\_binary =}\NormalTok{ il\_exposure }\SpecialCharTok{\textgreater{}}\NormalTok{ med\_il,      }\CommentTok{\# binary: high vs low IL exposure}
    \AttributeTok{treated\_cont   =}\NormalTok{ il\_exposure                }\CommentTok{\# continuous: raw IL routing share}
\NormalTok{  )}

\FunctionTok{cat}\NormalTok{(}\StringTok{"}\SpecialCharTok{\textbackslash{}n}\StringTok{Treated (high IL exposure):"}\NormalTok{, }\FunctionTok{sum}\NormalTok{(il\_exposure}\SpecialCharTok{$}\NormalTok{treated\_binary, }\AttributeTok{na.rm =} \ConstantTok{TRUE}\NormalTok{),}
    \StringTok{"municipalities}\SpecialCharTok{\textbackslash{}n}\StringTok{"}\NormalTok{)}
\end{Highlighting}
\end{Shaded}

\begin{Shaded}
\begin{Highlighting}[]
\NormalTok{il\_exposure }\SpecialCharTok{|\textgreater{}}
  \FunctionTok{filter}\NormalTok{(il\_exposure }\SpecialCharTok{\textgreater{}} \DecValTok{0}\NormalTok{) }\SpecialCharTok{|\textgreater{}}    \CommentTok{\# exclude municipalities with zero IL connection}
  \FunctionTok{ggplot}\NormalTok{(}\FunctionTok{aes}\NormalTok{(}\AttributeTok{x =}\NormalTok{ il\_exposure)) }\SpecialCharTok{+}
  \FunctionTok{geom\_histogram}\NormalTok{(}\AttributeTok{bins =} \DecValTok{60}\NormalTok{, }\AttributeTok{fill =} \StringTok{"darkgreen"}\NormalTok{, }\AttributeTok{alpha =} \FloatTok{0.75}\NormalTok{) }\SpecialCharTok{+}
  \FunctionTok{geom\_vline}\NormalTok{(}\AttributeTok{xintercept =}\NormalTok{ med\_il, }\AttributeTok{color =} \StringTok{"firebrick"}\NormalTok{,}
             \AttributeTok{linetype =} \StringTok{"dashed"}\NormalTok{, }\AttributeTok{linewidth =} \FloatTok{0.8}\NormalTok{) }\SpecialCharTok{+}
  \FunctionTok{annotate}\NormalTok{(}\StringTok{"text"}\NormalTok{, }\AttributeTok{x =}\NormalTok{ med\_il }\SpecialCharTok{+} \FloatTok{0.01}\NormalTok{, }\AttributeTok{y =} \ConstantTok{Inf}\NormalTok{, }\AttributeTok{vjust =} \FloatTok{1.5}\NormalTok{,}
           \AttributeTok{label =} \StringTok{"Median split"}\NormalTok{, }\AttributeTok{color =} \StringTok{"firebrick"}\NormalTok{, }\AttributeTok{size =} \DecValTok{3}\NormalTok{) }\SpecialCharTok{+}
  \FunctionTok{labs}\NormalTok{(}\AttributeTok{x =} \StringTok{"IL Routing Probability (row{-}normalised exposure)"}\NormalTok{,}
       \AttributeTok{y =} \StringTok{"Count"}\NormalTok{,}
       \AttributeTok{title =} \StringTok{"Distribution of Illinois Exposure Across MX Municipalities"}\NormalTok{) }\SpecialCharTok{+}
  \FunctionTok{theme\_minimal}\NormalTok{(}\AttributeTok{base\_size =} \DecValTok{9}\NormalTok{)}
\end{Highlighting}
\end{Shaded}

\subsection{Bridge Municipality
Identification}\label{bridge-municipality-identification}

Bridge municipalities are those that connect Illinois to California
through the network --- they send meaningful shares of migrants to both
states. We identify them using two criteria:

\begin{enumerate}
\def\labelenumi{\arabic{enumi}.}
\tightlist
\item
  \textbf{Network criterion:} top decile of betweenness centrality (from
  Section 2)
\item
  \textbf{Corridor criterion:} positive IL exposure AND positive CA
  exposure (both \textgreater{} 5\%)
\end{enumerate}

\begin{Shaded}
\begin{Highlighting}[]
\CommentTok{\# Identify bridge municipalities using both criteria}
\NormalTok{bridge\_munis }\OtherTok{\textless{}{-}}\NormalTok{ il\_exposure }\SpecialCharTok{|\textgreater{}}
  \FunctionTok{left\_join}\NormalTok{(}
\NormalTok{    exposure\_matrix }\SpecialCharTok{|\textgreater{}} \FunctionTok{select}\NormalTok{(mx\_municipality, }\AttributeTok{ca\_exposure =}\NormalTok{ California),}
    \AttributeTok{by =} \StringTok{"mx\_municipality"}
\NormalTok{  ) }\SpecialCharTok{|\textgreater{}}
  \FunctionTok{mutate}\NormalTok{(}
    \CommentTok{\# Corridor criterion: meaningful routing to BOTH IL and CA}
    \AttributeTok{corridor\_bridge =}\NormalTok{ il\_exposure }\SpecialCharTok{\textgreater{}} \FloatTok{0.05} \SpecialCharTok{\&}\NormalTok{ ca\_exposure }\SpecialCharTok{\textgreater{}} \FloatTok{0.05}\NormalTok{,}
    \CommentTok{\# Network criterion: top{-}decile betweenness (already in is\_bridge)}
    \CommentTok{\# Combined: either criterion}
    \AttributeTok{is\_bridge\_any =}\NormalTok{ is\_bridge }\SpecialCharTok{|}\NormalTok{ corridor\_bridge}
\NormalTok{  )}

\CommentTok{\# Tabulate}
\FunctionTok{cat}\NormalTok{(}\StringTok{"Bridge municipalities (network criterion):"}\NormalTok{, }\FunctionTok{sum}\NormalTok{(bridge\_munis}\SpecialCharTok{$}\NormalTok{is\_bridge,   }\AttributeTok{na.rm =} \ConstantTok{TRUE}\NormalTok{), }\StringTok{"}\SpecialCharTok{\textbackslash{}n}\StringTok{"}\NormalTok{)}
\FunctionTok{cat}\NormalTok{(}\StringTok{"Bridge municipalities (corridor criterion):"}\NormalTok{, }\FunctionTok{sum}\NormalTok{(bridge\_munis}\SpecialCharTok{$}\NormalTok{corridor\_bridge, }\AttributeTok{na.rm =} \ConstantTok{TRUE}\NormalTok{), }\StringTok{"}\SpecialCharTok{\textbackslash{}n}\StringTok{"}\NormalTok{)}
\FunctionTok{cat}\NormalTok{(}\StringTok{"Bridge municipalities (either criterion):  "}\NormalTok{, }\FunctionTok{sum}\NormalTok{(bridge\_munis}\SpecialCharTok{$}\NormalTok{is\_bridge\_any,   }\AttributeTok{na.rm =} \ConstantTok{TRUE}\NormalTok{), }\StringTok{"}\SpecialCharTok{\textbackslash{}n}\StringTok{"}\NormalTok{)}

\CommentTok{\# Show the top bridge municipalities}
\NormalTok{bridge\_munis }\SpecialCharTok{|\textgreater{}}
  \FunctionTok{filter}\NormalTok{(is\_bridge\_any) }\SpecialCharTok{|\textgreater{}}
  \FunctionTok{select}\NormalTok{(mx\_state, mx\_municipality, il\_exposure, ca\_exposure, betweenness, eigenvector) }\SpecialCharTok{|\textgreater{}}
  \FunctionTok{arrange}\NormalTok{(}\FunctionTok{desc}\NormalTok{(betweenness)) }\SpecialCharTok{|\textgreater{}}
  \FunctionTok{head}\NormalTok{(}\DecValTok{15}\NormalTok{) }\SpecialCharTok{|\textgreater{}}
\NormalTok{  knitr}\SpecialCharTok{::}\FunctionTok{kable}\NormalTok{(}\AttributeTok{digits =} \DecValTok{4}\NormalTok{,}
               \AttributeTok{caption =} \StringTok{"Top bridge municipalities by betweenness centrality."}\NormalTok{,}
               \AttributeTok{booktabs =} \ConstantTok{TRUE}\NormalTok{)}
\end{Highlighting}
\end{Shaded}

\subsection{Visualising the Impasse in Migration Time
Series}\label{visualising-the-impasse-in-migration-time-series}

Before any formal estimation, we inspect the migration time series for
bridge municipalities. The visual pattern is the primary face-validity
check: if the impasse had an effect, we expect Illinois-bound shares to
dip relative to California and New York during 2015--2017 for
municipalities with high IL exposure.

\begin{Shaded}
\begin{Highlighting}[]
\CommentTok{\# Use the long{-}format joint panel (already constructed in the data pipeline)}
\CommentTok{\# Filter to the top 6 bridge municipalities by betweenness}
\NormalTok{top6\_bridges }\OtherTok{\textless{}{-}}\NormalTok{ bridge\_munis }\SpecialCharTok{|\textgreater{}}
  \FunctionTok{filter}\NormalTok{(is\_bridge\_any) }\SpecialCharTok{|\textgreater{}}
  \FunctionTok{slice\_max}\NormalTok{(betweenness, }\AttributeTok{n =} \DecValTok{6}\NormalTok{) }\SpecialCharTok{|\textgreater{}}
  \FunctionTok{pull}\NormalTok{(mx\_municipality)}

\NormalTok{state\_colors }\OtherTok{\textless{}{-}} \FunctionTok{c}\NormalTok{(}\AttributeTok{california =} \StringTok{"steelblue"}\NormalTok{, }\AttributeTok{illinois =} \StringTok{"darkgreen"}\NormalTok{, }\AttributeTok{new\_york =} \StringTok{"firebrick"}\NormalTok{)}
\NormalTok{state\_labels }\OtherTok{\textless{}{-}} \FunctionTok{c}\NormalTok{(}\AttributeTok{california =} \StringTok{"California"}\NormalTok{, }\AttributeTok{illinois =} \StringTok{"Illinois"}\NormalTok{,  }\AttributeTok{new\_york =} \StringTok{"New York"}\NormalTok{)}

\NormalTok{long\_joint\_df }\SpecialCharTok{|\textgreater{}}
  \FunctionTok{filter}\NormalTok{(}\StringTok{\textasciigrave{}}\AttributeTok{origin muni}\StringTok{\textasciigrave{}} \SpecialCharTok{\%in\%}\NormalTok{ top6\_bridges, weight }\SpecialCharTok{\textgreater{}} \DecValTok{0}\NormalTok{) }\SpecialCharTok{|\textgreater{}}
  \FunctionTok{ggplot}\NormalTok{(}\FunctionTok{aes}\NormalTok{(}\AttributeTok{x =}\NormalTok{ year, }\AttributeTok{y =}\NormalTok{ weight, }\AttributeTok{color =}\NormalTok{ state)) }\SpecialCharTok{+}
  \FunctionTok{geom\_line}\NormalTok{(}\AttributeTok{linewidth =} \FloatTok{0.8}\NormalTok{) }\SpecialCharTok{+}
  \FunctionTok{geom\_point}\NormalTok{(}\AttributeTok{size =} \FloatTok{1.5}\NormalTok{) }\SpecialCharTok{+}
  \FunctionTok{annotate}\NormalTok{(}\StringTok{"rect"}\NormalTok{,}
           \AttributeTok{xmin =} \FloatTok{2015.5}\NormalTok{, }\AttributeTok{xmax =} \FloatTok{2017.5}\NormalTok{, }\AttributeTok{ymin =} \SpecialCharTok{{-}}\ConstantTok{Inf}\NormalTok{, }\AttributeTok{ymax =} \ConstantTok{Inf}\NormalTok{,}
           \AttributeTok{alpha =} \FloatTok{0.12}\NormalTok{, }\AttributeTok{fill =} \StringTok{"firebrick"}\NormalTok{) }\SpecialCharTok{+}
  \FunctionTok{geom\_vline}\NormalTok{(}\AttributeTok{xintercept =} \FunctionTok{c}\NormalTok{(}\FloatTok{2015.5}\NormalTok{, }\FloatTok{2017.5}\NormalTok{),}
             \AttributeTok{color =} \StringTok{"orange"}\NormalTok{, }\AttributeTok{linetype =} \StringTok{"dashed"}\NormalTok{, }\AttributeTok{linewidth =} \FloatTok{0.5}\NormalTok{) }\SpecialCharTok{+}
  \FunctionTok{scale\_color\_manual}\NormalTok{(}\AttributeTok{values =}\NormalTok{ state\_colors, }\AttributeTok{labels =}\NormalTok{ state\_labels) }\SpecialCharTok{+}
  \FunctionTok{scale\_x\_continuous}\NormalTok{(}\AttributeTok{breaks =} \DecValTok{2010}\SpecialCharTok{:}\DecValTok{2024}\NormalTok{, }\AttributeTok{labels =} \DecValTok{10}\SpecialCharTok{:}\DecValTok{24}\NormalTok{) }\SpecialCharTok{+}
  \FunctionTok{facet\_wrap}\NormalTok{(}\SpecialCharTok{\textasciitilde{}}\StringTok{\textasciigrave{}}\AttributeTok{origin muni}\StringTok{\textasciigrave{}}\NormalTok{, }\AttributeTok{scales =} \StringTok{"free\_y"}\NormalTok{, }\AttributeTok{ncol =} \DecValTok{3}\NormalTok{) }\SpecialCharTok{+}
  \FunctionTok{labs}\NormalTok{(}\AttributeTok{title    =} \StringTok{"Bridge Municipality Migration Shares (2010{-}2024)"}\NormalTok{,}
       \AttributeTok{subtitle =} \StringTok{"Red shading = IL Budget Impasse (2015{-}2017) | orange dashes = shock bounds"}\NormalTok{,}
       \AttributeTok{x =} \StringTok{"Year"}\NormalTok{, }\AttributeTok{y =} \StringTok{"Share of migrants"}\NormalTok{, }\AttributeTok{color =} \ConstantTok{NULL}\NormalTok{) }\SpecialCharTok{+}
  \FunctionTok{theme\_minimal}\NormalTok{(}\AttributeTok{base\_size =} \DecValTok{9}\NormalTok{) }\SpecialCharTok{+}
  \FunctionTok{theme}\NormalTok{(}\AttributeTok{legend.position =} \StringTok{"bottom"}\NormalTok{,}
        \AttributeTok{axis.text.x =} \FunctionTok{element\_text}\NormalTok{(}\AttributeTok{size =} \DecValTok{7}\NormalTok{))}
\end{Highlighting}
\end{Shaded}

\begin{center}\rule{0.5\linewidth}{0.5pt}\end{center}

\section{Difference-in-Differences
Estimation}\label{difference-in-differences-estimation}

\subsection{Setup and Panel
Construction}\label{setup-and-panel-construction}

Our DiD design compares the change in outcomes for \textbf{treated}
municipalities (high IL exposure) against \textbf{control}
municipalities (low IL exposure) before and after the Illinois Budget
Impasse. The identifying assumption is that in the absence of the shock,
treated and control municipalities would have followed parallel trends
in the outcome variable.

The formal DiD estimating equation is:

\[Y_{it} = \alpha_i + \gamma_t + \beta \cdot (\text{Treated}_i \times \text{Post}_t) + \varepsilon_{it}\]

where \(\alpha_i\) are municipality fixed effects (absorbing all
time-invariant municipality characteristics), \(\gamma_t\) are year
fixed effects (absorbing common shocks in all years), and \(\beta\) is
the average treatment effect of the Illinois Budget Impasse on outcome
\(Y\) for treated municipalities. We use \texttt{feols()} from the
\texttt{fixest} package, which implements the Mundlak-Chamberlain
estimator with clustered standard errors.

We study two outcome variables:

\begin{enumerate}
\def\labelenumi{\arabic{enumi}.}
\tightlist
\item
  \textbf{Direct effect:} \(Y_{it} = \hat{R}_{it}^{IL}\) ---
  municipality \(i\)'s imputed remittance inflow \emph{via the Illinois
  corridor} in year \(t\). A negative \(\hat{\beta}\) confirms that the
  impasse depressed migration-mediated remittance flows for IL-exposed
  municipalities.
\item
  \textbf{Second-order propagation:} \(Y_{it} = W_{i,CA,t}\) ---
  municipality \(i\)'s migration share to California in year \(t\). A
  positive \(\hat{\beta}\) for bridge municipalities would mean that as
  Illinois-bound flows fell, some fraction was redirected toward
  California, confirming cross-state propagation.
\end{enumerate}

\begin{Shaded}
\begin{Highlighting}[]
\CommentTok{\# Aggregate R\_hat to annual, restricting to the Illinois corridor}
\CommentTok{\# (R\_hat\_df has one row per municipality and one column per quarter;}
\CommentTok{\#  we sum the quarterly Illinois{-}corridor component, not total R\_hat)}
\NormalTok{R\_hat\_il\_annual }\OtherTok{\textless{}{-}}\NormalTok{ R\_hat\_df }\SpecialCharTok{|\textgreater{}}
  \FunctionTok{pivot\_longer}\NormalTok{(}
    \AttributeTok{cols      =} \SpecialCharTok{{-}}\FunctionTok{c}\NormalTok{(mx\_state, mx\_municipality),}
    \AttributeTok{names\_to  =} \StringTok{"date"}\NormalTok{,}
    \AttributeTok{values\_to =} \StringTok{"r\_hat"}
\NormalTok{  ) }\SpecialCharTok{|\textgreater{}}
  \FunctionTok{mutate}\NormalTok{(}\AttributeTok{year =}\NormalTok{ lubridate}\SpecialCharTok{::}\FunctionTok{year}\NormalTok{(}\FunctionTok{as.Date}\NormalTok{(date))) }\SpecialCharTok{|\textgreater{}}
  \FunctionTok{group\_by}\NormalTok{(mx\_state, mx\_municipality, year) }\SpecialCharTok{|\textgreater{}}
  \FunctionTok{summarise}\NormalTok{(}\AttributeTok{r\_hat\_total =} \FunctionTok{sum}\NormalTok{(r\_hat, }\AttributeTok{na.rm =} \ConstantTok{TRUE}\NormalTok{), }\AttributeTok{.groups =} \StringTok{"drop"}\NormalTok{)}

\CommentTok{\# Merge with treatment variables from il\_exposure}
\NormalTok{panel\_df }\OtherTok{\textless{}{-}}\NormalTok{ R\_hat\_il\_annual }\SpecialCharTok{|\textgreater{}}
  \FunctionTok{left\_join}\NormalTok{(}
\NormalTok{    il\_exposure }\SpecialCharTok{|\textgreater{}}
      \FunctionTok{select}\NormalTok{(mx\_state, mx\_municipality,}
\NormalTok{             il\_exposure, treated\_binary, treated\_cont,}
\NormalTok{             betweenness, eigenvector, eigen\_std, is\_bridge),}
    \AttributeTok{by =} \FunctionTok{c}\NormalTok{(}\StringTok{"mx\_state"}\NormalTok{, }\StringTok{"mx\_municipality"}\NormalTok{)}
\NormalTok{  ) }\SpecialCharTok{|\textgreater{}}
  \CommentTok{\# Define the post{-}treatment indicator}
  \CommentTok{\# Post = 1 for the impasse years (2015, 2016, 2017)}
  \CommentTok{\# Pre  = 0 for 2010{-}2014}
  \CommentTok{\# We exclude 2018{-}2024 from the primary DiD to keep the window clean}
  \FunctionTok{filter}\NormalTok{(year }\SpecialCharTok{\%in\%} \DecValTok{2010}\SpecialCharTok{:}\DecValTok{2019}\NormalTok{) }\SpecialCharTok{|\textgreater{}}
  \FunctionTok{mutate}\NormalTok{(}
    \AttributeTok{post         =} \FunctionTok{as.integer}\NormalTok{(year }\SpecialCharTok{\%in\%} \DecValTok{2015}\SpecialCharTok{:}\DecValTok{2017}\NormalTok{),   }\CommentTok{\# 1 during impasse, 0 before}
    \AttributeTok{treat\_x\_post =}\NormalTok{ treated\_binary }\SpecialCharTok{*}\NormalTok{ post,              }\CommentTok{\# DiD interaction term}
    \AttributeTok{log\_r\_hat    =} \FunctionTok{log}\NormalTok{(r\_hat\_total }\SpecialCharTok{+} \DecValTok{1}\NormalTok{),               }\CommentTok{\# log{-}transform for stability}
    \CommentTok{\# Create a clean panel ID}
    \AttributeTok{muni\_id      =} \FunctionTok{paste}\NormalTok{(mx\_state, mx\_municipality, }\AttributeTok{sep =} \StringTok{"\_\_"}\NormalTok{)}
\NormalTok{  )}

\FunctionTok{cat}\NormalTok{(}\StringTok{"Panel dimensions:"}\NormalTok{, }\FunctionTok{nrow}\NormalTok{(panel\_df), }\StringTok{"rows,"}\NormalTok{,}
    \FunctionTok{n\_distinct}\NormalTok{(panel\_df}\SpecialCharTok{$}\NormalTok{muni\_id), }\StringTok{"municipalities,"}\NormalTok{,}
    \FunctionTok{n\_distinct}\NormalTok{(panel\_df}\SpecialCharTok{$}\NormalTok{year), }\StringTok{"years}\SpecialCharTok{\textbackslash{}n}\StringTok{"}\NormalTok{)}
\end{Highlighting}
\end{Shaded}

\subsection{Direct Effect: Illinois Exposure -\textgreater{} Imputed
Remittance
Drop}\label{direct-effect-illinois-exposure---imputed-remittance-drop}

The first test is whether high-IL-exposure municipalities experienced a
relative decline in their imputed remittance inflows during the impasse
period.

\begin{Shaded}
\begin{Highlighting}[]
\CommentTok{\# Two{-}way fixed effects DiD using feols() from fixest.}
\CommentTok{\# | muni\_id + year means: absorb municipality and year fixed effects.}
\CommentTok{\# cluster = \textasciitilde{}muni\_id clusters standard errors at the municipality level.}
\CommentTok{\# This is the standard Bertrand{-}Duflo{-}Mullainathan correction for serial correlation.}

\CommentTok{\# Model 1: binary treatment, log R\_hat outcome}
\NormalTok{m1 }\OtherTok{\textless{}{-}} \FunctionTok{feols}\NormalTok{(}
\NormalTok{  log\_r\_hat }\SpecialCharTok{\textasciitilde{}}\NormalTok{ treat\_x\_post }\SpecialCharTok{|}\NormalTok{ muni\_id }\SpecialCharTok{+}\NormalTok{ year,}
  \AttributeTok{data    =}\NormalTok{ panel\_df,}
  \AttributeTok{cluster =} \SpecialCharTok{\textasciitilde{}}\NormalTok{muni\_id}
\NormalTok{)}

\CommentTok{\# Model 2: continuous treatment (raw IL exposure share)}
\NormalTok{m2 }\OtherTok{\textless{}{-}} \FunctionTok{feols}\NormalTok{(}
\NormalTok{  log\_r\_hat }\SpecialCharTok{\textasciitilde{}}\NormalTok{ il\_exposure}\SpecialCharTok{:}\NormalTok{post }\SpecialCharTok{|}\NormalTok{ muni\_id }\SpecialCharTok{+}\NormalTok{ year,}
  \AttributeTok{data    =}\NormalTok{ panel\_df,}
  \AttributeTok{cluster =} \SpecialCharTok{\textasciitilde{}}\NormalTok{muni\_id}
\NormalTok{)}

\CommentTok{\# Model 3: eigenvector centrality as continuous treatment intensity}
\NormalTok{m3 }\OtherTok{\textless{}{-}} \FunctionTok{feols}\NormalTok{(}
\NormalTok{  log\_r\_hat }\SpecialCharTok{\textasciitilde{}}\NormalTok{ eigen\_std}\SpecialCharTok{:}\NormalTok{post }\SpecialCharTok{|}\NormalTok{ muni\_id }\SpecialCharTok{+}\NormalTok{ year,}
  \AttributeTok{data    =}\NormalTok{ panel\_df,}
  \AttributeTok{cluster =} \SpecialCharTok{\textasciitilde{}}\NormalTok{muni\_id}
\NormalTok{)}

\CommentTok{\# Regression table}
\FunctionTok{modelsummary}\NormalTok{(}
  \FunctionTok{list}\NormalTok{(}\StringTok{"Binary (IL \textgreater{} median)"} \OtherTok{=}\NormalTok{ m1,}
       \StringTok{"Continuous (IL share)"} \OtherTok{=}\NormalTok{ m2,}
       \StringTok{"Eigenvector centrality"} \OtherTok{=}\NormalTok{ m3),}
  \AttributeTok{stars      =} \FunctionTok{c}\NormalTok{(}\StringTok{"*"} \OtherTok{=} \FloatTok{0.10}\NormalTok{, }\StringTok{"**"} \OtherTok{=} \FloatTok{0.05}\NormalTok{, }\StringTok{"***"} \OtherTok{=} \FloatTok{0.01}\NormalTok{),}
  \AttributeTok{coef\_rename =} \FunctionTok{c}\NormalTok{(}
    \StringTok{"treat\_x\_post"}     \OtherTok{=} \StringTok{"Treated x Post"}\NormalTok{,}
    \StringTok{"il\_exposure:post"} \OtherTok{=} \StringTok{"IL Exposure x Post"}\NormalTok{,}
    \StringTok{"eigen\_std:post"}   \OtherTok{=} \StringTok{"Eigenvector Centrality x Post"}
\NormalTok{  ),}
  \AttributeTok{gof\_omit   =} \StringTok{"AIC|BIC|RMSE|R2 Within|Std"}\NormalTok{,}
  \AttributeTok{title      =} \StringTok{"Direct Effect: Illinois Budget Impasse on Imputed Remittances"}\NormalTok{,}
  \AttributeTok{notes      =} \StringTok{"Two{-}way FE (municipality + year). SE clustered by municipality. Outcome: log(R{-}hat + 1)."}\NormalTok{,}
  \AttributeTok{output     =} \StringTok{"kableExtra"}
\NormalTok{)}
\end{Highlighting}
\end{Shaded}

\begin{Shaded}
\begin{Highlighting}[]
\CommentTok{\# Event study: interact treatment with year dummies to trace out the treatment effect}
\CommentTok{\# year by year. Flat pre{-}trends validate the parallel{-}trends assumption.}
\CommentTok{\# Reference year: 2014 (one year before the shock)}
\NormalTok{panel\_df }\OtherTok{\textless{}{-}}\NormalTok{ panel\_df }\SpecialCharTok{|\textgreater{}}
  \FunctionTok{mutate}\NormalTok{(}\AttributeTok{year\_f =} \FunctionTok{relevel}\NormalTok{(}\FunctionTok{factor}\NormalTok{(year), }\AttributeTok{ref =} \StringTok{"2014"}\NormalTok{))}

\NormalTok{m\_event }\OtherTok{\textless{}{-}} \FunctionTok{feols}\NormalTok{(}
\NormalTok{  log\_r\_hat }\SpecialCharTok{\textasciitilde{}}\NormalTok{ treated\_binary}\SpecialCharTok{:}\NormalTok{year\_f }\SpecialCharTok{|}\NormalTok{ muni\_id }\SpecialCharTok{+}\NormalTok{ year,}
  \AttributeTok{data    =}\NormalTok{ panel\_df,}
  \AttributeTok{cluster =} \SpecialCharTok{\textasciitilde{}}\NormalTok{muni\_id}
\NormalTok{)}

\CommentTok{\# Extract coefficients and confidence intervals for the event study plot}
\NormalTok{event\_coefs }\OtherTok{\textless{}{-}}\NormalTok{ broom}\SpecialCharTok{::}\FunctionTok{tidy}\NormalTok{(m\_event, }\AttributeTok{conf.int =} \ConstantTok{TRUE}\NormalTok{) }\SpecialCharTok{|\textgreater{}}
  \FunctionTok{filter}\NormalTok{(}\FunctionTok{str\_detect}\NormalTok{(term, }\StringTok{"year\_f"}\NormalTok{)) }\SpecialCharTok{|\textgreater{}}
  \FunctionTok{mutate}\NormalTok{(}
    \AttributeTok{year =} \FunctionTok{as.integer}\NormalTok{(}\FunctionTok{str\_extract}\NormalTok{(term, }\StringTok{"}\SpecialCharTok{\textbackslash{}\textbackslash{}}\StringTok{d\{4\}"}\NormalTok{))}
\NormalTok{  ) }\SpecialCharTok{|\textgreater{}}
  \CommentTok{\# Bind a zero row for the reference year 2014}
  \FunctionTok{bind\_rows}\NormalTok{(}\FunctionTok{tibble}\NormalTok{(}\AttributeTok{year =} \DecValTok{2014}\NormalTok{, }\AttributeTok{estimate =} \DecValTok{0}\NormalTok{, }\AttributeTok{conf.low =} \DecValTok{0}\NormalTok{, }\AttributeTok{conf.high =} \DecValTok{0}\NormalTok{))}

\FunctionTok{ggplot}\NormalTok{(event\_coefs, }\FunctionTok{aes}\NormalTok{(}\AttributeTok{x =}\NormalTok{ year, }\AttributeTok{y =}\NormalTok{ estimate)) }\SpecialCharTok{+}
  \FunctionTok{geom\_hline}\NormalTok{(}\AttributeTok{yintercept =} \DecValTok{0}\NormalTok{, }\AttributeTok{color =} \StringTok{"grey50"}\NormalTok{, }\AttributeTok{linewidth =} \FloatTok{0.5}\NormalTok{) }\SpecialCharTok{+}
  \FunctionTok{geom\_ribbon}\NormalTok{(}\FunctionTok{aes}\NormalTok{(}\AttributeTok{ymin =}\NormalTok{ conf.low, }\AttributeTok{ymax =}\NormalTok{ conf.high), }\AttributeTok{alpha =} \FloatTok{0.15}\NormalTok{, }\AttributeTok{fill =} \StringTok{"steelblue"}\NormalTok{) }\SpecialCharTok{+}
  \FunctionTok{geom\_line}\NormalTok{(}\AttributeTok{color =} \StringTok{"steelblue"}\NormalTok{, }\AttributeTok{linewidth =} \FloatTok{0.8}\NormalTok{) }\SpecialCharTok{+}
  \FunctionTok{geom\_point}\NormalTok{(}\AttributeTok{color =} \StringTok{"steelblue"}\NormalTok{, }\AttributeTok{size =} \DecValTok{2}\NormalTok{) }\SpecialCharTok{+}
  \FunctionTok{annotate}\NormalTok{(}\StringTok{"rect"}\NormalTok{, }\AttributeTok{xmin =} \FloatTok{2015.5}\NormalTok{, }\AttributeTok{xmax =} \FloatTok{2017.5}\NormalTok{, }\AttributeTok{ymin =} \SpecialCharTok{{-}}\ConstantTok{Inf}\NormalTok{, }\AttributeTok{ymax =} \ConstantTok{Inf}\NormalTok{,}
           \AttributeTok{alpha =} \FloatTok{0.08}\NormalTok{, }\AttributeTok{fill =} \StringTok{"firebrick"}\NormalTok{) }\SpecialCharTok{+}
  \FunctionTok{geom\_vline}\NormalTok{(}\AttributeTok{xintercept =} \FunctionTok{c}\NormalTok{(}\FloatTok{2015.5}\NormalTok{, }\FloatTok{2017.5}\NormalTok{),}
             \AttributeTok{color =} \StringTok{"firebrick"}\NormalTok{, }\AttributeTok{linetype =} \StringTok{"dashed"}\NormalTok{, }\AttributeTok{linewidth =} \FloatTok{0.5}\NormalTok{) }\SpecialCharTok{+}
  \FunctionTok{scale\_x\_continuous}\NormalTok{(}\AttributeTok{breaks =} \DecValTok{2010}\SpecialCharTok{:}\DecValTok{2019}\NormalTok{) }\SpecialCharTok{+}
  \FunctionTok{labs}\NormalTok{(}\AttributeTok{x =} \StringTok{"Year"}\NormalTok{, }\AttributeTok{y =} \StringTok{"Coefficient (relative to 2014)"}\NormalTok{,}
       \AttributeTok{title =} \StringTok{"Event Study: Treated vs Control Municipalities"}\NormalTok{,}
       \AttributeTok{subtitle =} \StringTok{"Pre{-}trend flat {-}\textgreater{} parallel trends. Drop during red window = treatment effect."}\NormalTok{) }\SpecialCharTok{+}
  \FunctionTok{theme\_minimal}\NormalTok{(}\AttributeTok{base\_size =} \DecValTok{9}\NormalTok{)}
\end{Highlighting}
\end{Shaded}

\subsection{Second-Order Propagation: Bridge Municipalities
-\textgreater{}
California}\label{second-order-propagation-bridge-municipalities---california}

The direct effect confirms that Illinois-exposed municipalities were
hurt by the impasse. The propagation hypothesis goes further: did those
municipalities \emph{redirect} migrant flows toward California, creating
an indirect second-order effect on the CA corridor?

We test this by restricting to \textbf{bridge municipalities} (those
with meaningful ties to both IL and CA) and estimating a DiD on their
California migration shares. The treatment indicator here is high IL
exposure \emph{conditional on being a bridge municipality}: we want to
know whether, among municipalities that connect to both states, those
more IL-dependent showed a larger CA share \emph{increase} during the
impasse.

\begin{Shaded}
\begin{Highlighting}[]
\CommentTok{\# Build the California{-}outcome panel for bridge municipalities only.}
\CommentTok{\# Outcome: migration share to California (from long\_joint\_df)}
\CommentTok{\# Treatment: IL exposure (already defined)}
\CommentTok{\# We restrict to bridge municipalities (is\_bridge\_any = TRUE)}

\NormalTok{panel\_ca }\OtherTok{\textless{}{-}}\NormalTok{ long\_joint\_df }\SpecialCharTok{|\textgreater{}}
  \FunctionTok{filter}\NormalTok{(state }\SpecialCharTok{==} \StringTok{"california"}\NormalTok{) }\SpecialCharTok{|\textgreater{}}                      \CommentTok{\# California outcome}
  \FunctionTok{rename}\NormalTok{(}\StringTok{\textasciigrave{}}\AttributeTok{origin muni}\StringTok{\textasciigrave{}} \OtherTok{=} \StringTok{\textasciigrave{}}\AttributeTok{origin muni}\StringTok{\textasciigrave{}}\NormalTok{,}
         \AttributeTok{ca\_share =}\NormalTok{ weight) }\SpecialCharTok{|\textgreater{}}
  \FunctionTok{left\_join}\NormalTok{(}
\NormalTok{    bridge\_munis }\SpecialCharTok{|\textgreater{}}
      \FunctionTok{select}\NormalTok{(mx\_municipality, il\_exposure, is\_bridge\_any,}
\NormalTok{             treated\_binary, eigenvector, eigen\_std),}
    \AttributeTok{by =} \FunctionTok{c}\NormalTok{(}\StringTok{"origin muni"} \OtherTok{=} \StringTok{"mx\_municipality"}\NormalTok{)}
\NormalTok{  ) }\SpecialCharTok{|\textgreater{}}
  \FunctionTok{filter}\NormalTok{(is\_bridge\_any, year }\SpecialCharTok{\%in\%} \DecValTok{2010}\SpecialCharTok{:}\DecValTok{2019}\NormalTok{) }\SpecialCharTok{|\textgreater{}}         \CommentTok{\# bridge munis only}
  \FunctionTok{mutate}\NormalTok{(}
    \AttributeTok{post          =} \FunctionTok{as.integer}\NormalTok{(year }\SpecialCharTok{\%in\%} \DecValTok{2015}\SpecialCharTok{:}\DecValTok{2017}\NormalTok{),}
    \AttributeTok{treat\_x\_post  =}\NormalTok{ treated\_binary }\SpecialCharTok{*}\NormalTok{ post,}
    \AttributeTok{muni\_id       =} \StringTok{\textasciigrave{}}\AttributeTok{origin muni}\StringTok{\textasciigrave{}}\NormalTok{,}
    \AttributeTok{log\_ca\_share  =} \FunctionTok{log}\NormalTok{(ca\_share }\SpecialCharTok{+} \FloatTok{1e{-}6}\NormalTok{)               }\CommentTok{\# log{-}transform (add small constant)}
\NormalTok{  )}

\FunctionTok{cat}\NormalTok{(}\StringTok{"Propagation panel: "}\NormalTok{, }\FunctionTok{nrow}\NormalTok{(panel\_ca), }\StringTok{"rows,"}\NormalTok{,}
    \FunctionTok{n\_distinct}\NormalTok{(panel\_ca}\SpecialCharTok{$}\NormalTok{muni\_id), }\StringTok{"bridge municipalities}\SpecialCharTok{\textbackslash{}n}\StringTok{"}\NormalTok{)}
\end{Highlighting}
\end{Shaded}

\begin{Shaded}
\begin{Highlighting}[]
\CommentTok{\# Second{-}order propagation DiD.}
\CommentTok{\# H0: no propagation {-}\textgreater{} treat\_x\_post coefficient is zero}
\CommentTok{\# H1: propagation    {-}\textgreater{} treat\_x\_post \textgreater{} 0}
\CommentTok{\#   (more IL{-}exposed bridges redirect more flow to CA during the impasse)}

\CommentTok{\# Model 1: binary treatment}
\NormalTok{p1 }\OtherTok{\textless{}{-}} \FunctionTok{feols}\NormalTok{(}
\NormalTok{  log\_ca\_share }\SpecialCharTok{\textasciitilde{}}\NormalTok{ treat\_x\_post }\SpecialCharTok{|}\NormalTok{ muni\_id }\SpecialCharTok{+}\NormalTok{ year,}
  \AttributeTok{data    =}\NormalTok{ panel\_ca,}
  \AttributeTok{cluster =} \SpecialCharTok{\textasciitilde{}}\NormalTok{muni\_id}
\NormalTok{)}

\CommentTok{\# Model 2: continuous IL exposure}
\NormalTok{p2 }\OtherTok{\textless{}{-}} \FunctionTok{feols}\NormalTok{(}
\NormalTok{  log\_ca\_share }\SpecialCharTok{\textasciitilde{}}\NormalTok{ il\_exposure}\SpecialCharTok{:}\NormalTok{post }\SpecialCharTok{|}\NormalTok{ muni\_id }\SpecialCharTok{+}\NormalTok{ year,}
  \AttributeTok{data    =}\NormalTok{ panel\_ca,}
  \AttributeTok{cluster =} \SpecialCharTok{\textasciitilde{}}\NormalTok{muni\_id}
\NormalTok{)}

\CommentTok{\# Model 3: eigenvector centrality as treatment intensity}
\NormalTok{p3 }\OtherTok{\textless{}{-}} \FunctionTok{feols}\NormalTok{(}
\NormalTok{  log\_ca\_share }\SpecialCharTok{\textasciitilde{}}\NormalTok{ eigen\_std}\SpecialCharTok{:}\NormalTok{post }\SpecialCharTok{|}\NormalTok{ muni\_id }\SpecialCharTok{+}\NormalTok{ year,}
  \AttributeTok{data    =}\NormalTok{ panel\_ca,}
  \AttributeTok{cluster =} \SpecialCharTok{\textasciitilde{}}\NormalTok{muni\_id}
\NormalTok{)}

\FunctionTok{modelsummary}\NormalTok{(}
  \FunctionTok{list}\NormalTok{(}\StringTok{"Binary (IL \textgreater{} median)"} \OtherTok{=}\NormalTok{ p1,}
       \StringTok{"Continuous (IL share)"} \OtherTok{=}\NormalTok{ p2,}
       \StringTok{"Eigenvector centrality"} \OtherTok{=}\NormalTok{ p3),}
  \AttributeTok{stars      =} \FunctionTok{c}\NormalTok{(}\StringTok{"*"} \OtherTok{=} \FloatTok{0.10}\NormalTok{, }\StringTok{"**"} \OtherTok{=} \FloatTok{0.05}\NormalTok{, }\StringTok{"***"} \OtherTok{=} \FloatTok{0.01}\NormalTok{),}
  \AttributeTok{coef\_rename =} \FunctionTok{c}\NormalTok{(}
    \StringTok{"treat\_x\_post"}     \OtherTok{=} \StringTok{"Treated x Post"}\NormalTok{,}
    \StringTok{"il\_exposure:post"} \OtherTok{=} \StringTok{"IL Exposure x Post"}\NormalTok{,}
    \StringTok{"eigen\_std:post"}   \OtherTok{=} \StringTok{"Eigenvector Centrality x Post"}
\NormalTok{  ),}
  \AttributeTok{gof\_omit   =} \StringTok{"AIC|BIC|RMSE|R2 Within|Std"}\NormalTok{,}
  \AttributeTok{title      =} \StringTok{"Second{-}Order Propagation: Bridge Municipality CA Shares During IL Impasse"}\NormalTok{,}
  \AttributeTok{notes      =} \StringTok{"Bridge municipalities only. Two{-}way FE (muni + year). SE clustered by municipality. Outcome: log(CA migration share + eps)."}\NormalTok{,}
  \AttributeTok{output     =} \StringTok{"kableExtra"}
\NormalTok{)}
\end{Highlighting}
\end{Shaded}

\subsection{Placebo: Non-Bridge
Municipalities}\label{placebo-non-bridge-municipalities}

If the propagation result is real, it should be \emph{specific to bridge
municipalities}. Municipalities with zero or negligible CA ties should
show no CA-share increase during the impasse, even if they are highly
IL-exposed. This is a within-design placebo check.

\begin{Shaded}
\begin{Highlighting}[]
\CommentTok{\# Placebo group: high{-}IL municipalities that are NOT bridges (no meaningful CA ties)}
\NormalTok{panel\_placebo }\OtherTok{\textless{}{-}}\NormalTok{ long\_joint\_df }\SpecialCharTok{|\textgreater{}}
  \FunctionTok{filter}\NormalTok{(state }\SpecialCharTok{==} \StringTok{"california"}\NormalTok{, year }\SpecialCharTok{\%in\%} \DecValTok{2010}\SpecialCharTok{:}\DecValTok{2019}\NormalTok{) }\SpecialCharTok{|\textgreater{}}
  \FunctionTok{rename}\NormalTok{(}\AttributeTok{ca\_share =}\NormalTok{ weight) }\SpecialCharTok{|\textgreater{}}
  \FunctionTok{left\_join}\NormalTok{(}
\NormalTok{    bridge\_munis }\SpecialCharTok{|\textgreater{}}
      \FunctionTok{select}\NormalTok{(mx\_municipality, il\_exposure, is\_bridge\_any, treated\_binary),}
    \AttributeTok{by =} \FunctionTok{c}\NormalTok{(}\StringTok{"origin muni"} \OtherTok{=} \StringTok{"mx\_municipality"}\NormalTok{)}
\NormalTok{  ) }\SpecialCharTok{|\textgreater{}}
  \FunctionTok{filter}\NormalTok{(}\SpecialCharTok{!}\NormalTok{is\_bridge\_any) }\SpecialCharTok{|\textgreater{}}             \CommentTok{\# NON{-}bridge municipalities only}
  \FunctionTok{mutate}\NormalTok{(}
    \AttributeTok{post         =} \FunctionTok{as.integer}\NormalTok{(year }\SpecialCharTok{\%in\%} \DecValTok{2015}\SpecialCharTok{:}\DecValTok{2017}\NormalTok{),}
    \AttributeTok{treat\_x\_post =}\NormalTok{ treated\_binary }\SpecialCharTok{*}\NormalTok{ post,}
    \AttributeTok{muni\_id      =} \StringTok{\textasciigrave{}}\AttributeTok{origin muni}\StringTok{\textasciigrave{}}\NormalTok{,}
    \AttributeTok{log\_ca\_share =} \FunctionTok{log}\NormalTok{(ca\_share }\SpecialCharTok{+} \FloatTok{1e{-}6}\NormalTok{)}
\NormalTok{  )}

\NormalTok{p\_placebo }\OtherTok{\textless{}{-}} \FunctionTok{feols}\NormalTok{(}
\NormalTok{  log\_ca\_share }\SpecialCharTok{\textasciitilde{}}\NormalTok{ treat\_x\_post }\SpecialCharTok{|}\NormalTok{ muni\_id }\SpecialCharTok{+}\NormalTok{ year,}
  \AttributeTok{data    =}\NormalTok{ panel\_placebo,}
  \AttributeTok{cluster =} \SpecialCharTok{\textasciitilde{}}\NormalTok{muni\_id}
\NormalTok{)}

\FunctionTok{modelsummary}\NormalTok{(}
  \FunctionTok{list}\NormalTok{(}\StringTok{"Bridge municipalities"} \OtherTok{=}\NormalTok{ p1,}
       \StringTok{"Non{-}bridge (placebo)"}  \OtherTok{=}\NormalTok{ p\_placebo),}
  \AttributeTok{stars      =} \FunctionTok{c}\NormalTok{(}\StringTok{"*"} \OtherTok{=} \FloatTok{0.10}\NormalTok{, }\StringTok{"**"} \OtherTok{=} \FloatTok{0.05}\NormalTok{, }\StringTok{"***"} \OtherTok{=} \FloatTok{0.01}\NormalTok{),}
  \AttributeTok{coef\_rename =} \FunctionTok{c}\NormalTok{(}\StringTok{"treat\_x\_post"} \OtherTok{=} \StringTok{"Treated x Post"}\NormalTok{),}
  \AttributeTok{gof\_omit   =} \StringTok{"AIC|BIC|RMSE|R2 Within|Std"}\NormalTok{,}
  \AttributeTok{title      =} \StringTok{"Propagation vs Placebo: Bridge vs Non{-}Bridge Municipalities"}\NormalTok{,}
  \AttributeTok{notes      =} \StringTok{"If the network story is correct, the treatment effect should be significant for bridges and near{-}zero for non{-}bridges."}\NormalTok{,}
  \AttributeTok{output     =} \StringTok{"kableExtra"}
\NormalTok{)}
\end{Highlighting}
\end{Shaded}

\begin{center}\rule{0.5\linewidth}{0.5pt}\end{center}

\section{Discussion}\label{discussion}

\subsection{The Cascade
Interpretation}\label{the-cascade-interpretation}

The results, taken together, tell a coherent cascade story. The Illinois
Budget Impasse constituted a localised, negative labor-market shock that
operated through the migration network in three sequential stages.

\textbf{Stage 1 -- Direct hit.} Municipalities whose migrants were
concentrated in Illinois experienced a measurable decline in imputed
remittance inflows during the impasse years. The DiD coefficient on the
direct effect (Section 4.1) captures this first-order impact: every
percentage-point increase in pre-impasse IL routing probability is
associated with an approximately \([INSERT ESTIMATE]\%\) relative
decline in imputed remittances during 2015--2017. The event study
confirms that the effect is concentrated precisely in the impasse
window, with flat pre-trends from 2010--2014.

\textbf{Stage 2 -- Network transmission.} Bridge municipalities ---
those sitting at high betweenness positions in the bipartite graph and
maintaining simultaneous IL and CA corridors --- were the transmission
nodes. The community detection analysis in Section 2 established that
Illinois and California belong to the same migration community in the
\(M^\top M\) graph, meaning they recruit from overlapping sending pools.
When Illinois migration became less attractive, migrants from bridge
municipalities had a ready alternative corridor: California. The
second-order propagation DiD (Section 4.2) tests this directly and finds
\([INSERT ESTIMATE]\) for bridge municipalities versus a near-zero
coefficient for the non-bridge placebo group.

\textbf{Stage 3 -- Asymmetric response.} Not all states in Illinois's
community would be equally affected. States that are closer substitutes
in network space (higher \(M^\top M\) overlap with Illinois) should
absorb more of the redirected flow. Quantifying this substitution
elasticity --- \(\Delta W_{i,CA} / \Delta W_{i,IL}\) across bridge
municipalities --- is a natural extension of this work.

\subsection{The TRUST Act as a Positive Shock Robustness
Check}\label{the-trust-act-as-a-positive-shock-robustness-check}

The Illinois TRUST Act (signed August 28, 2017) provides an important
robustness check with the opposite sign. The TRUST Act reduced the
deportation risk for undocumented migrants in Illinois, which should
have made Illinois \emph{more} attractive relative to other states. If
our network mechanism is correct, we should observe:

\begin{enumerate}
\def\labelenumi{\arabic{enumi}.}
\tightlist
\item
  A positive direct effect: IL-exposed municipalities show
  \emph{increasing} migration shares to Illinois post-2017
\item
  A negative second-order effect: bridge municipalities show
  \emph{decreasing} CA shares as some flow returns to Illinois
\end{enumerate}

\begin{Shaded}
\begin{Highlighting}[]
\CommentTok{\# TRUST Act DiD: same design, different post window}
\CommentTok{\# Post = 2018 and 2019 (full years after TRUST Act signed Aug 2017)}
\CommentTok{\# Pre  = 2010{-}2014 (same pre{-}period as before for a clean comparison)}
\NormalTok{panel\_trust }\OtherTok{\textless{}{-}}\NormalTok{ panel\_df }\SpecialCharTok{|\textgreater{}}
  \FunctionTok{filter}\NormalTok{(year }\SpecialCharTok{\%in\%} \FunctionTok{c}\NormalTok{(}\DecValTok{2010}\SpecialCharTok{:}\DecValTok{2014}\NormalTok{, }\DecValTok{2018}\SpecialCharTok{:}\DecValTok{2019}\NormalTok{)) }\SpecialCharTok{|\textgreater{}}  \CommentTok{\# exclude the impasse years}
  \FunctionTok{mutate}\NormalTok{(}
    \AttributeTok{post\_trust   =} \FunctionTok{as.integer}\NormalTok{(year }\SpecialCharTok{\%in\%} \DecValTok{2018}\SpecialCharTok{:}\DecValTok{2019}\NormalTok{),}
    \AttributeTok{treat\_x\_post =}\NormalTok{ treated\_binary }\SpecialCharTok{*}\NormalTok{ post\_trust}
\NormalTok{  )}

\NormalTok{m\_trust }\OtherTok{\textless{}{-}} \FunctionTok{feols}\NormalTok{(}
\NormalTok{  log\_r\_hat }\SpecialCharTok{\textasciitilde{}}\NormalTok{ treat\_x\_post }\SpecialCharTok{|}\NormalTok{ muni\_id }\SpecialCharTok{+}\NormalTok{ year,}
  \AttributeTok{data    =}\NormalTok{ panel\_trust,}
  \AttributeTok{cluster =} \SpecialCharTok{\textasciitilde{}}\NormalTok{muni\_id}
\NormalTok{)}

\FunctionTok{modelsummary}\NormalTok{(}
  \FunctionTok{list}\NormalTok{(}\StringTok{"Impasse (negative shock)"} \OtherTok{=}\NormalTok{ m1,}
       \StringTok{"TRUST Act (positive shock)"} \OtherTok{=}\NormalTok{ m\_trust),}
  \AttributeTok{stars      =} \FunctionTok{c}\NormalTok{(}\StringTok{"*"} \OtherTok{=} \FloatTok{0.10}\NormalTok{, }\StringTok{"**"} \OtherTok{=} \FloatTok{0.05}\NormalTok{, }\StringTok{"***"} \OtherTok{=} \FloatTok{0.01}\NormalTok{),}
  \AttributeTok{coef\_rename =} \FunctionTok{c}\NormalTok{(}\StringTok{"treat\_x\_post"} \OtherTok{=} \StringTok{"Treated x Post"}\NormalTok{),}
  \AttributeTok{gof\_omit   =} \StringTok{"AIC|BIC|RMSE|R2 Within|Std"}\NormalTok{,}
  \AttributeTok{title      =} \StringTok{"Robustness: TRUST Act as Opposite{-}Sign Shock"}\NormalTok{,}
  \AttributeTok{notes      =} \StringTok{"If the network mechanism is symmetric, the TRUST Act coefficient should be positive and significant, opposite to the impasse coefficient."}\NormalTok{,}
  \AttributeTok{output     =} \StringTok{"kableExtra"}
\NormalTok{)}
\end{Highlighting}
\end{Shaded}

A positive and significant coefficient on the TRUST Act interaction,
combined with a negative coefficient for the impasse, would constitute
strong evidence that the migration network is a genuine bidirectional
transmission channel --- not just a one-time ratchet that only responds
to negative shocks.

\subsection{Policy Implications}\label{policy-implications}

\textbf{For migrant-receiving states.} The results suggest that US
states are not isolated labor markets for Mexican migrants. A state that
creates conditions hostile to migrants (as the budget impasse
effectively did) does not simply retain fewer migrants --- it pushes
them toward network-adjacent states, generating positive spillovers for
those states. Conversely, states that adopt migrant-friendly policies
(like the TRUST Act) may attract inflows partly diverted from other
corridors. This has implications for how we think about the political
economy of immigration policy: states that free-ride on other states'
enforcement face positive network externalities.

\textbf{For Mexican sending communities.} Bridge municipalities are
particularly exposed to cross-state volatility: their remittance income
can fluctuate based on fiscal crises in states they are only partially
tied to. Identifying these structurally vulnerable communities in
advance --- using precisely the betweenness centrality measure developed
in Section 2 --- could guide targeted remittance stabilisation programs
or financial inclusion policies in Mexico.

\textbf{For identification strategy design.} Our analysis demonstrates
that network structure can be used \emph{predictively} to motivate DiD
designs before looking at the outcome data. The community detection
result in Section 2 predicted, on purely structural grounds, that
Illinois shocks should propagate to California and not to states in
other communities. Researchers designing DiD studies in networked
settings should consider running community detection on the relevant
adjacency matrix as a pre-registration step: it constrains which control
groups are valid and which secondary outcomes are worth testing,
reducing the risk of data mining.

\subsection{Limitations and
Extensions}\label{limitations-and-extensions}

Several important caveats apply. First, MCAS data captures only migrants
who registered at Mexican consulates, which is a subset of the total
migrant population. To the extent that registration propensity
correlates with IL exposure (e.g., more documented workers are
proportionally more affected by the impasse), our estimates will be
attenuated. Second, we use imputed remittances \(\hat{R}\) as the
outcome rather than DS1 actual remittances, both because DS1 has
incomplete municipality coverage and because it is measured with error.
The \(r = 0.79\) validation in the data pipeline suggests \(\hat{R}\) is
a credible proxy, but any measurement error in the outcome is a source
of downward bias. Third, the parallel-trends assumption requires that
treated and control municipalities would have moved together in the
absence of the impasse --- a maintained assumption we can only partially
validate through pre-trends in the event study.

Future work should explore: (i) heterogeneity in the propagation
coefficient by community overlap (\(M^\top M_{j,CA}\)) to test whether
closer substitutes receive larger redirected flows; (ii) the full
general equilibrium response, including whether bridge municipalities
also show changes in flows to New York; and (iii) a structural network
model that can separate the migration-rerouting channel from a direct
remittance-income channel operating through migrants who stay in
Illinois but reduce remittances.




\end{document}
